\chapter{Structure of rigid local systems}
\label{chap:5}
\setcounter{section}{-1}

\section{Cohomological rigidity}\label{sec:5.0}

\sourcepara{5.0.1}
Let \(\ell\) be a prime number, \(k\) an algebraically closed field of
characteristic \(\neq \ell\). On \(\mathbb{A}^1\) over \(k\), denote by
\(\kappa:\mathbb{A}^1\to \mathbb{P}^1\) the inclusion. Let \(K\) in
\(\mathrm{D}^b_c(\mathbb{A}^1,\overline{\mathbb{Q}}_\ell)\) be a perverse
irreducible sheaf which is nonpunctual, i.e., \(K\) is \(j_*\mathcal{F}[1]\)
for a nonempty open set \(j:U\to\mathbb{A}^1\), and a lisse irreducible
\(\overline{\mathbb{Q}}_\ell\)-sheaf \(\mathcal{F}\) on \(U\). In
\hyperref[par:3.0.1]{(3.0.1)}, we defined the index of rigidity of \(K\),
\(\operatorname{rig}(K)\), to be the integer
\[
  \operatorname{rig}(K):=
  \chi(\mathbb{P}^1,\kappa_*j_*\underline{\operatorname{End}}(\mathcal{F})).
\]
We say that \(K\), or \(\mathcal{F}\), is cohomologically rigid, if
\(\operatorname{rig}(K)=2\).

\sourceheading[thm:5.0.2]{Theorem 5.0.2}
Let \(\ell\) be a prime number, \(k\) an algebraically closed field of
characteristic \(\neq\ell\). In \(\mathbb{A}^1\) over \(k\), let
\(j:U\to\mathbb{A}^1\) be a nonempty open set, \(\mathcal{F}\) a lisse
irreducible \(\overline{\mathbb{Q}}_\ell\)-sheaf on \(U\). Suppose that
\(\mathcal{F}\) is cohomologically rigid. Let \(\mathcal{G}\) be any lisse
\(\overline{\mathbb{Q}}_\ell\)-sheaf on \(U\) which is locally isomorphic to
\(\mathcal{F}\) in the sense that for each point \(s\) of \(\mathbb{P}^1-U\),
there exists an isomorphism between the representations \(\mathcal{F}(s)\) and
\(\mathcal{G}(s)\) of the inertia group \(I(s)\) afforded by \(\mathcal{F}\)
and \(\mathcal{G}\) respectively. Then there exists an isomorphism
\(\mathcal{F}\cong\mathcal{G}\) of lisse \(\overline{\mathbb{Q}}_\ell\)-sheaf
on \(U\).

\begin{proof}
This is a trivial modification of the easy half of
\hyperref[thm:1.1.2]{Theorem~1.1.2}. Denote by
\(\kappa:\mathbb{A}^1\to\mathbb{P}^1\) the inclusion. Since \(\mathcal{F}\)
and \(\mathcal{G}\) are locally isomorphic, so are the two local systems
\(\underline{\operatorname{End}}(\mathcal{F})\) and
\(\underline{\operatorname{Hom}}(\mathcal{F},\mathcal{G})\). So from the shape
of the Euler--Poincaré formula, which depends only on ranks and on local data
at the ramification points and at infinity, we see that
\[
  2=\chi(\mathbb{P}^1,\kappa_*j_*\underline{\operatorname{End}}(\mathcal{F}))
  =\chi(\mathbb{P}^1,\kappa_*j_*\underline{\operatorname{Hom}}(\mathcal{F},\mathcal{G})).
\]
Once we have the inequality
\(\chi(\mathbb{P}^1,\kappa_*j_*\underline{\operatorname{Hom}}(\mathcal{F},\mathcal{G}))\geq 2\),
we conclude exactly as we did in \hyperref[thm:1.1.2]{Theorem~1.1.2}.
\end{proof}

\section{The category \texorpdfstring{\(\mathcal{T}_\ell\)}{T ell}, and the
functors \texorpdfstring{\(\operatorname{MC}_\chi\)}{MC chi} and
\texorpdfstring{\(\operatorname{MT}_{\mathcal{L}}\)}{MT L}}
\label{sec:5.1}

\sourcepara{5.1.1}
In this section, we work on \(\mathbb{A}^1\) over an algebraically closed
field \(k\). We fix a prime number \(\ell\neq\operatorname{char}(k)\). We are
interested in the full subcategory \(\mathcal{T}_\ell\) of the category of
constructible \(\overline{\mathbb{Q}}_\ell\)-sheaves \(\mathcal{F}\) on
\(\mathbb{A}^1\) consisting of those which satisfy the following three
conditions:
\begin{enumerate}[label={\(\mathcal{T}\arabic*\)}]
\item \(\mathcal{F}\) is an irreducible middle extension: there exists a dense
open set \(j:U\to\mathbb{A}^1\) such that \(j^*\mathcal{F}\) is lisse and
irreducible on \(U\), and such that \(\mathcal{F}\cong j_*j^*\mathcal{F}\).

\item \(\mathcal{F}\) is tame: \(j^*\mathcal{F}\) is tamely ramified at every
point of \(\mathbb{P}^1-U\).

\item \(\mathcal{F}\) has at least two finite singularities: there are at least
two distinct points of \(\mathbb{A}^1\) at which \(\mathcal{F}\) fails to be
lisse.
\end{enumerate}

\sourceheading[lem:5.1.2]{Lemma 5.1.2}
In the situation \hyperref[par:5.0.1]{(5.0.1)}, suppose \(\mathcal{F}\)
satisfies conditions \(\mathcal{T}1\) and \(\mathcal{T}2\) above. If
\(\mathcal{F}\) has generic rank \(\geq 2\), then \(\mathcal{F}\) satisfies
\(\mathcal{T}3\).

\begin{proof}
Indeed, if \(\mathcal{F}\) is tame and irreducible with at most one finite
singularity, say \(\alpha\), then \(j^*\mathcal{F}\) is an irreducible
representation of the abelian group
\(\pi_1^{\mathrm{tame}}(\mathbb{A}^1-\{\alpha\})\), so has rank \(\leq 1\).
\end{proof}

\sourceheading[lem:5.1.3]{Lemma 5.1.3}
In the situation \hyperref[par:5.0.1]{(5.0.1)}, the perverse irreducibles
\(K\) in \(\mathrm{D}^b_c(\mathbb{A}^1,\overline{\mathbb{Q}}_\ell)\) which are
tame, and of type \(2d\)) in the sense of
\hyperref[thm:3.3.3]{Theorem~3.3.3} are precisely those of the form
\(\mathcal{F}[1]\), with \(\mathcal{F}\) in \(\mathcal{T}_\ell\).

\begin{proof}
A nonpunctual perverse irreducible \(K\) is precisely an \(\mathcal{F}[1]\) for
\(\mathcal{F}\) satisfying condition \(\mathcal{T}1\) above (and we recover
\(\mathcal{F}\) as \(\mathcal{H}^{-1}(K)\)). Such a perverse irreducible \(K\)
is tame if and only if \(\mathcal{F}\) is tame. Requiring in addition that
\(K\) have property \(\mathcal{P}\) forces \(\mathcal{F}\) to be nonconstant
(and not to be an \(\mathcal{L}_\psi\), but that case was already eliminated by
the tameness requirement). Requiring \(K\) to be of type \(2d\)) eliminates the
possibility that \(\mathcal{F}\) be any
\(j_*\mathcal{L}_{\chi(x-\alpha)}\). But these last two eliminations, of
\(\overline{\mathbb{Q}}_\ell\) and of any
\(j_*\mathcal{L}_{\chi(x-\alpha)}\), amount precisely to condition
\(\mathcal{T}3\) above. Indeed, of tame irreducible middle extensions, those
with at most one finite singularity are either \(\overline{\mathbb{Q}}_\ell\)
(if there is no finite singularity, because
\(\pi_1^{\mathrm{tame}}(\mathbb{A}^1)=0\)) or they are
\(j_*\mathcal{L}_{\chi(x-\alpha)}\) for \(\alpha\) the unique finite
singularity (because \(\pi_1^{\mathrm{tame}}(\mathbb{G}_m)\) is abelian).
\end{proof}

\sourcepara{5.1.4}
For any nontrivial continuous character
\[
  \chi:\pi_1^{\mathrm{tame}}(\mathbb{G}_{m,k})\to
  \overline{\mathbb{Q}}_\ell^\times,
\]
we denote by \(\mathcal{L}_\chi\) the corresponding lisse sheaf of rank one on
\(\mathbb{G}_m\), and by \(j_*\mathcal{L}_\chi\) its direct image to
\(\mathbb{A}^1\). We may restate \hyperref[thm:4.3.11]{Theorem~4.3.11} in
terms of \(\mathcal{T}_\ell\) as follows. Given a finite subset \(D\) in
\(\mathbb{A}^1(k)\), and an \(\mathcal{F}\) in \(\mathcal{T}_\ell\) which is
lisse on \(\mathbb{A}^1-D\), the middle convolution
\[
  \mathcal{F}[1]*_{\mathrm{mid}+}j_*\mathcal{L}_\chi[1]
\]
is of the form \(\mathcal{G}[1]\), with \(\mathcal{G}\) in
\(\mathcal{T}_\ell\), and \(\mathcal{G}\) lisse on the same
\(\mathbb{A}^1-D\). We can recover \(\mathcal{F}\) from \(\mathcal{G}\) by the
inversion formula
\[
  \mathcal{F}[1]=\mathcal{G}[1]*_{\mathrm{mid}+}j_*\mathcal{L}_{\bar{\chi}}[1].
\]
Moreover, the local monodromies of \(\mathcal{F}\) and of \(\mathcal{G}\) are
related by \hyperref[cor:3.3.6]{Corollary~3.3.6} and
\hyperref[cor:3.3.7]{Corollary~3.3.7}. According to
\hyperref[thm:4.3.10]{Theorem~4.3.10}, the indices of rigidity of
\(\mathcal{F}\) and \(\mathcal{G}\) (strictly speaking, of the perverse objects
\(\mathcal{F}[1]\) and of \(\mathcal{G}[1]\)) are equal:
\[
  \operatorname{rig}(\mathcal{F})=\operatorname{rig}(\mathcal{G}).
\]

\sourcepara{5.1.5}
We define, for \(\chi\) a nontrivial character as above, the functor
\[
  \operatorname{MC}_\chi:\mathcal{T}_\ell\to\mathcal{T}_\ell
\]
by
\[
  \operatorname{MC}_\chi(\mathcal{F}):=\mathcal{G}:=
  (\mathcal{F}[1]*_{\mathrm{mid}+}j_*\mathcal{L}_\chi[1])[-1].
\]
These functors have the composition laws
\[
\begin{aligned}
  \operatorname{MC}_\chi\circ\operatorname{MC}_\rho
  &=\operatorname{MC}_\rho\circ\operatorname{MC}_\chi
  =\operatorname{MC}_{\chi\rho}\quad\text{if \(\chi\rho\) is nontrivial},\\
  \operatorname{MC}_\chi\circ\operatorname{MC}_{\bar{\chi}}&=\mathrm{id}.
\end{aligned}
\]

\sourcepara{5.1.6}
Here is a relatively concrete description of \(\operatorname{MC}_\chi(\mathcal{F})\).
Let \(\mathcal{F}\) in \(\mathcal{T}_\ell\). Denote by \(D\) the set of finite
singularities of \(\mathcal{F}\), and by
\[
  j:\mathbb{A}^1-D\to\mathbb{A}^1
\]
the inclusion. Consider the projection
\[
  \operatorname{pr}_2:\mathbb{A}^2\to\mathbb{A}^1,\quad (x,t)\mapsto t.
\]
On \(\mathbb{A}^2\), we have the sheaf
\(\mathcal{F}_x\otimes\mathcal{L}_{\chi(t-x)}\) (extended by zero across the
diagonal \(t=x\)). Because \(\mathcal{F}\) is lisse on \(\mathbb{A}^1-D\), and
in \(\mathcal{T}_\ell\), the higher direct images
\(\mathrm{R}^i\operatorname{pr}_{2!}(\mathcal{F}_x\otimes\mathcal{L}_{\chi(t-x)})\)
and
\(\mathrm{R}^i\operatorname{pr}_{2*}(\mathcal{F}_x\otimes\mathcal{L}_{\chi(t-x)})\)
are both lisse on \(\mathbb{A}^1-D\), and both vanish for \(i\neq 1\).
Moreover, both are of formation compatible with arbitrary change of base on
\(\mathbb{A}^1-D\) (by \hyperref[cor:2.8.5]{Corollary~2.8.5}). Thus we may
form the lisse sheaf on \(\mathbb{A}^1-D\) which is
\[
  \operatorname{Image}\bigl(
  j^*\mathrm{R}^1\operatorname{pr}_{2!}(\mathcal{F}_x\otimes\mathcal{L}_{\chi(t-x)})
  \to
  j^*\mathrm{R}^1\operatorname{pr}_{2*}(\mathcal{F}_x\otimes\mathcal{L}_{\chi(t-x)})
  \bigr).
\]
It is tautological that this is none other than
\(j^*\operatorname{MC}_\chi(\mathcal{F})\). But
\(\operatorname{MC}_\chi(\mathcal{F})\) is a middle extension, so we get the

\sourceheading[recipe:5.1.7]{Explicit Recipes 5.1.7}
In the situation \hyperref[par:5.1.1]{(5.1.1)}, we have:
\begin{enumerate}[label=(\arabic*)]
\item For \(\mathcal{F}\) in \(\mathcal{T}_\ell\), and
\(j:\mathbb{A}^1-D\to\mathbb{A}^1\) the inclusion of a dense open set on which
\(j^*\mathcal{F}\) is lisse,
\[
\begin{aligned}
  \operatorname{MC}_\chi(\mathcal{F})
  &=j_*j^*\operatorname{MC}_\chi(\mathcal{F})\\
  &=j_*\operatorname{Image}\bigl(
  j^*\mathrm{R}^1\operatorname{pr}_{2!}(\mathcal{F}_x\otimes\mathcal{L}_{\chi(t-x)})
  \to
  j^*\mathrm{R}^1\operatorname{pr}_{2*}(\mathcal{F}_x\otimes\mathcal{L}_{\chi(t-x)})
  \bigr).
\end{aligned}
\]

\item At any geometric point \(\tau\) in \(\mathbb{A}^1-D\), the stalk at
\(\tau\) of \(\operatorname{MC}_\chi(\mathcal{F})\) is given by
\[
  (\operatorname{MC}_\chi(\mathcal{F}))_\tau
  =\operatorname{Image}\bigl(
  \mathrm{H}^1_c(\mathbb{A}^1,\mathcal{F}\otimes\mathcal{L}_{\chi(\tau-x)})
  \to
  \mathrm{H}^1(\mathbb{A}^1,\mathcal{F}\otimes\mathcal{L}_{\chi(\tau-x)})
  \bigr).
\]

\item In terms of the inclusion \(\kappa:\mathbb{A}^1\to\mathbb{P}^1\), this
image is just the ``parabolic'' group
\(\mathrm{H}^1(\mathbb{P}^1,\kappa_*(\mathcal{F}\otimes\mathcal{L}_{\chi(\tau-x)}))\),
so
\[
  (\operatorname{MC}_\chi(\mathcal{F}))_\tau
  =\mathrm{H}^1(\mathbb{P}^1,\kappa_*(\mathcal{F}\otimes\mathcal{L}_{\chi(\tau-x)})),
  \quad\text{for \(\tau\) in \(\mathbb{A}^1-D\)}.
\]

\item In terms of the inclusion
\(\kappa_1:\mathbb{A}^1-D-\{\tau\}\to\mathbb{P}^1\), we may rewrite this as
\[
\begin{aligned}
  (\operatorname{MC}_\chi(\mathcal{F}))_\tau
  &=
  \mathrm{H}^1(\mathbb{P}^1,\kappa_{1*}((\mathcal{F}\otimes
  \mathcal{L}_{\chi(\tau-x)})\vert\mathbb{A}^1-D-\{\tau\}))\\
  &=
  \operatorname{Image}\text{ of the ``forget supports'' map}\\
  &\quad
  \mathrm{H}^1_c(\mathbb{A}^1-D-\{\tau\},
  \mathcal{F}\otimes\mathcal{L}_{\chi(\tau-x)})
  \to
  \mathrm{H}^1(\mathbb{A}^1-D-\{\tau\},
  \mathcal{F}\otimes\mathcal{L}_{\chi(\tau-x)}).
\end{aligned}
\]
\end{enumerate}

\sourcepara{5.1.8}
Let \(\mathcal{L}\) be a middle extension sheaf on \(\mathbb{A}^1\) which is
generically of rank one, and which is tame. We wish to define an operation of
``middle tensor product with \(\mathcal{L}\)'' for objects of
\(\mathcal{T}_\ell\) (compare \hyperref[sec:3.2]{\S3.2}). The naive idea is
this: given \(\mathcal{F}\) in \(\mathcal{T}_\ell\), pick
\(j:U\to\mathbb{A}^1\) a dense open on which both \(\mathcal{F}\) and
\(\mathcal{L}\) are lisse, and form
\[
  j_*((j^*\mathcal{F})\otimes(j^*\mathcal{L})).
\]
This sheaf, which is independent of the auxiliary choice of \(U\), visibly
satisfies both \(\mathcal{T}1\) and \(\mathcal{T}2\). As it has the same
generic rank as \(\mathcal{F}\), it will satisfy \(\mathcal{T}3\) if
\(\mathcal{F}\) has generic rank at least \(2\), thanks to
\hyperref[lem:5.1.2]{Lemma~5.1.2}.

\sourcepara{5.1.9}
However, for \(\mathcal{F}\) of generic rank one,
\(j_*((j^*\mathcal{F})\otimes(j^*\mathcal{L}))\) may very well fail to satisfy
\(\mathcal{T}3\): for example, it might be the constant sheaf. In order to deal
with this problem, we define
\[
  \mathcal{T}_{\ell,\operatorname{rk}\geq 2}:=
  \text{the full subcategory of \(\mathcal{T}_\ell\) consisting of those
  objects of generic rank \(\geq 2\)}.
\]
On this full subcategory \(\mathcal{T}_{\ell,\operatorname{rk}\geq 2}\), we can
define the above ``middle tensor product with \(\mathcal{L}\)'', which we denote
\[
\begin{aligned}
  \operatorname{MT}_{\mathcal{L}}:
  \mathcal{T}_{\ell,\operatorname{rk}\geq 2}\to
  \mathcal{T}_{\ell,\operatorname{rk}\geq 2},\\
  \operatorname{MT}_{\mathcal{L}}(\mathcal{F}):=
  j_*((j^*\mathcal{F})\otimes(j^*\mathcal{L})),
\end{aligned}
\]
for \(j:U\to\mathbb{A}^1\) any dense open on which both \(\mathcal{F}\) and
\(\mathcal{L}\) are lisse. It is clear from the definitions that
\(\operatorname{MT}_{\mathcal{L}}\) preserves index of rigidity:
\[
  \operatorname{rig}(\operatorname{MT}_{\mathcal{L}}(\mathcal{F}))
  =\operatorname{rig}(\mathcal{F}).
\]
If we are given a lisse, rank one \(\mathcal{L}\) on a dense open
\(j:U\to\mathbb{A}^1\), we will often write \(\operatorname{MT}_{\mathcal{L}}\)
instead of the literally correct \(\operatorname{MT}_{j_*\mathcal{L}}\).

\section{The main theorem on the structure of rigid local systems}\label{sec:5.2}

\sourcepara{5.2.0}
Given an everywhere tame lisse \(\overline{\mathbb{Q}}_\ell\)-sheaf on
\(\mathbb{A}^1-D\), we denote by \(\operatorname{LocMono}(\mathcal{F})\) the
subgroup of the group of all continuous characters
\[
  \chi:I(0)^{\mathrm{tame}}\cong
  \pi_1^{\mathrm{tame}}(\mathbb{G}_{m,k})\to
  \overline{\mathbb{Q}}_\ell^\times
\]
generated by those which ``occur'' in the local monodromies of \(\mathcal{F}\)
at all the points of \(\mathbb{P}^1\) (i.e., canonically identify the tame
inertia groups \(I(\alpha)^{\mathrm{tame}}\) at points \(\alpha\) in
\(\mathbb{P}^1\) to \(I(0)^{\mathrm{tame}}\) via automorphisms of
\(\mathbb{P}^1\) which carry \(\alpha\) to \(0\) (the resulting identifications
are independent of the choice)).

\sourceheading[thm:5.2.1]{Main Theorem 5.2.1}
In the situation \hyperref[par:5.1.1]{(5.1.1)}, suppose \(\mathcal{F}\) in
\(\mathcal{T}_\ell\) has generic rank \(r(\mathcal{F})\geq 2\), and is lisse
on \(\mathbb{A}^1-D\). Suppose further that \(\mathcal{F}\) is
cohomologically rigid. Then
\begin{enumerate}[label=(\arabic*)]
\item There exists a lisse, everywhere tame, rank one
\(\overline{\mathbb{Q}}_\ell\)-sheaf \(\mathcal{L}\) on \(\mathbb{A}^1-D\),
and a nontrivial character \(\chi\), such that the (necessarily
cohomologically rigid) object \(\mathcal{G}\) in \(\mathcal{T}_\ell\) defined
by
\[
  \mathcal{G}:=\operatorname{MC}_\chi\operatorname{MT}_{\mathcal{L}}(\mathcal{F})
\]
has strictly lower generic rank: \(r(\mathcal{G})<r(\mathcal{F})\).

\item In (1) above, we may choose both \(\mathcal{L}_\chi\) and
\(\mathcal{L}\) to have all of their local monodromies in the group
\(\operatorname{LocMono}(\mathcal{F})\), and, if we so choose them, then
\(\mathcal{G}:=\operatorname{MC}_\chi\operatorname{MT}_{\mathcal{L}}(\mathcal{F})\)
has \(\operatorname{LocMono}(\mathcal{G})\subset\operatorname{LocMono}(\mathcal{F})\).

\item If for some integer \(N\geq 1\) all the eigenvalues of all the local
monodromies of \(\mathcal{F}\) are \(N\)'th roots of unity, then in (1) above
we may choose \(\chi\) of order dividing \(N\), and we may choose
\(\mathcal{L}\) to have each of its local monodromies of order dividing \(N\),
and if we so choose them,
all the eigenvalues of all the local monodromies of \(\mathcal{G}\) will be
\(N\)'th roots of unity.
\end{enumerate}

\sourcepara{5.2.2}
\textit{Proof.} We will prove the theorem by giving an algorithm for how to
choose both \(\mathcal{L}\) and \(\chi\).

\sourcepara{5.2.2.1}
In order to eliminate any ambiguity in the algorithm, for each point \(\alpha\)
in \(\mathbb{P}^1(k)\), we will choose a total ordering of the underlying set
of the group of all continuous characters
\[
  \chi:I(\alpha)^{\mathrm{tame}}\to\overline{\mathbb{Q}}_\ell^\times.
\]
One way to do this is as follows. If we pick a topological generator, say
\(\gamma^{\mathrm{tame}}\), of the group \(I(\alpha)^{\mathrm{tame}}\) (which
is canonically the group
\(\prod_{\ell\neq\operatorname{char}(k)}\mathbb{Z}_\ell(1)\)), and a field
embedding of \(\overline{\mathbb{Q}}_\ell\) into \(\mathbb{C}\), then
``evaluation at \(\gamma^{\mathrm{tame}}\)'' embeds the continuous characters
into the group \(\mathbb{C}^\times\). So we are reduced to exhibiting a total
ordering of \(\mathbb{C}^\times\). The simplest way to do this is to use polar
coordinates, writing nonzero complex numbers as \(re^{i\vartheta}\) with \(r\)
and \(\vartheta\) real, \(r>0\) and \(0\leq \vartheta<2\pi\). Then we say
\(re^{i\vartheta}<r'e^{i\vartheta'}\) if either \(r<r'\) or if \(r=r'\) and
\(\vartheta<\vartheta'\).

\sourcepara{5.2.2.2}
\textbf{First Step.} We now explain how to choose \(\mathcal{L}\). Because
\(\mathcal{F}\) is lisse on \(\mathbb{A}^1-D\), and everywhere tame, at each
point \(\alpha\) in \(D\amalg\{\infty\}\), the local monodromy of
\(\mathcal{F}\) defines an \(r(\mathcal{F})\)-dimensional representation of
\(I(\alpha)^{\mathrm{tame}}\). We first describe a shorthand to describe this
representation of \(I(\alpha)^{\mathrm{tame}}\). We write it as a direct sum of
(character)\(\otimes\)(unipotent representation):
\[
\begin{aligned}
  \mathcal{F}\text{ as }I(\alpha)^{\mathrm{tame}}\text{-rep'n.}
  &=\bigoplus_\chi
  \mathcal{L}_{\chi(x-\alpha)}\otimes
  \operatorname{Unip}(\alpha,\chi,\mathcal{F}),\quad \alpha\text{ finite},\\
  \mathcal{F}\text{ as }I(\infty)^{\mathrm{tame}}\text{-rep'n.}
  &=\bigoplus_\chi
  \mathcal{L}_\chi\otimes\operatorname{Unip}(\infty,\chi,\mathcal{F}),
  \quad\text{for \(\alpha=\infty\)}.
\end{aligned}
\]
We write \(\operatorname{Unip}(\alpha,\chi,\mathcal{F})\) as a direct sum of
Jordan blocks, of dimensions \(\{n_i(\alpha,\chi,\mathcal{F})\}_i\), and then
pass to its ``dual partition'' (compare \hyperref[par:3.1.10]{(3.1.10)}), the
decreasing sequence of non-negative integers
\[
  e_1(\alpha,\chi,\mathcal{F})\geq e_2(\alpha,\chi,\mathcal{F})\geq\cdots
  \geq e_k(\alpha,\chi,\mathcal{F})=0\quad\text{for \(k\gg 0\)}
\]
defined by
\[
  e_j(\alpha,\chi,\mathcal{F}):=
  \text{the number of Jordan blocks in }
  \operatorname{Unip}(\alpha,\chi,\mathcal{F})
  \text{ whose dimension is \(\geq j\)}.
\]
Thus \(e_1(\alpha,\chi,\mathcal{F})\) is the multiplicity
\[
  \dim\operatorname{Hom}_{I(\alpha)}
  (\{\mathcal{L}_{\chi(x-\alpha)}\text{ if \(\alpha\) finite, }
  \mathcal{L}_\chi\text{ if \(\alpha=\infty\)}\},\mathcal{F}(\alpha))
\]
of \(\chi\) as simple eigenvalue in
\(\mathcal{F}(\alpha):=\mathcal{F}\) as \(I(\alpha)^{\mathrm{tame}}\)-representation
(and \(e_j(\alpha,\chi,\mathcal{F})=0\) for all \(j\) if the character \(\chi\)
does not occur in the local monodromy of \(\mathcal{F}\) at \(\alpha\)).

Given \(\mathcal{F}\) and a point \(\alpha\) in \(D\), pick a character
\(\chi\) whose \(e_1(\alpha,\chi,\mathcal{F})\) is maximal among all the
characters which occur in the local monodromy of \(\mathcal{F}\) at \(\alpha\).
There may not be a unique such \(\chi\), but in case of a tie we choose the one
which comes first in the chosen total ordering of the set of all characters of
\(I(\alpha)^{\mathrm{tame}}\). (The proof works just as well if whenever we
are confronted with a tie we make an arbitrary choice, rather than consult some
pre-chosen total ordering.) This character we denote \(\chi_{\alpha,\mathcal{F}}\).

\sourcepara{5.2.2.3}
We now have chosen, at each point \(\alpha\) in the finite set
\(D\subset\mathbb{A}^1\), a continuous character \(\chi_{\alpha,\mathcal{F}}\)
of \(I(\alpha)^{\mathrm{tame}}\). By the known structure of the prime-to-\(p\)
fundamental group of \(\mathbb{A}^1-D\) (the maximal prime-to-\(p\) quotient of
the profinite completion of the free group with one generator for each point
\(\alpha\) in \(D\)), there is a unique (up to isomorphism) lisse, everywhere
tame, rank one \(\overline{\mathbb{Q}}_\ell\)-sheaf \(\mathcal{L}\) on
\(\mathbb{A}^1-D\) with the property that, for every \(\alpha\) in \(D\),
\[
  \mathcal{L}\text{ as }I(\alpha)^{\mathrm{tame}}\text{-representation}
  =(\chi_{\alpha,\mathcal{F}})^{-1}.
\]
Concretely, \(\mathcal{L}\) is the tensor product, over \(\alpha\) in \(D\), of
the inverses of the translated Kummer sheaves
\(\mathcal{L}_{\chi_{\alpha,\mathcal{F}}(x-\alpha)}\).

\sourcepara{5.2.2.4}
With this choice of \(\mathcal{L}\), the sheaf
\(\operatorname{MT}_{\mathcal{L}}(\mathcal{F})\), which has the same generic
rank as \(\mathcal{F}\) and the same index of rigidity as \(\mathcal{F}\), has
the additional property that, at each \(\alpha\) in \(D\), the trivial
character \(\mathbf{1}\) has maximal ``simple multiplicity'': for each
\(\alpha\) in \(D\), we have
\[
  e_1(\alpha,\mathbf{1},\operatorname{MT}_{\mathcal{L}}(\mathcal{F}))
  \geq
  e_1(\alpha,\chi,\operatorname{MT}_{\mathcal{L}}(\mathcal{F}))
  \quad\text{for all \(\chi\)}.
\]

\sourcepara{5.2.2.5}
Notice that the local monodromies of \(\mathcal{L}\) were all drawn from the
group \(\operatorname{LocMono}(\mathcal{F})\), and that, consequently,
\(\operatorname{LocMono}(\operatorname{MT}_{\mathcal{L}}(\mathcal{F}))\) lies
in \(\operatorname{LocMono}(\mathcal{F})\).

\sourcepara{5.2.2.6}
\textbf{Second Step.} In order to complete the proof of the theorem, it
suffices to apply the following result to
\(\operatorname{MT}_{\mathcal{L}}(\mathcal{F})\).

\sourceheading[thm:5.2.3]{Main Theorem 5.2.3 (= \hyperref[thm:5.2.1]{\textcolor{blue}{5.2.1 bis}})}
In the situation \hyperref[par:5.1.1]{(5.1.1)}, suppose \(\mathcal{F}\) in
\(\mathcal{T}_\ell\) has generic rank \(r(\mathcal{F})\geq 2\), and is lisse
on \(\mathbb{A}^1-D\). Suppose that \(\mathcal{F}\) satisfies the following
condition:
\[
  (*)\quad
  \text{At every point \(\alpha\) in \(D\), we have }
  e_1(\alpha,\mathbf{1},\mathcal{F})\geq e_1(\alpha,\chi,\mathcal{F})
  \text{ for all \(\chi\)}.
\]
If \(\mathcal{F}\) is in addition cohomologically rigid, there exists a
nontrivial character \(\chi\) in \(\operatorname{LocMono}(\mathcal{F})\), such
that the (necessarily cohomologically rigid) object \(\mathcal{G}\) in
\(\mathcal{T}_\ell\) defined by
\[
  \mathcal{G}:=\operatorname{MC}_\chi(\mathcal{F})
\]
has strictly lower generic rank: \(r(\mathcal{G})<r(\mathcal{F})\). Moreover,
\(\operatorname{LocMono}(\mathcal{G})\) is contained in
\(\operatorname{LocMono}(\mathcal{F})\).

\sourcepara{5.2.4}
\textit{Proof.} The ``moreover'' follows from the explicit recipes of
\hyperref[cor:3.3.6]{Corollary~3.3.6}.

\sourcepara{5.2.4.1}
We choose the character \(\chi\) so as to maximize the dimension of the
\(I(\infty)^{\mathrm{tame}}\) invariants in
\(\mathcal{L}_\chi\otimes\mathcal{F}\). (If there is more than one \(\chi\)
which works, we choose the first of them in the chosen total ordering, though
any choice would work for the proof.)

\sourcepara{5.2.4.2}
We claim that any \(\chi\) which maximize the dimension of the
\(I(\infty)^{\mathrm{tame}}\) invariants in
\(\mathcal{L}_\chi\otimes\mathcal{F}\) is nontrivial. If not, we would have
\[
  e_1(\infty,\mathbf{1},\mathcal{F})\geq
  e_1(\infty,\chi,\mathcal{F})\quad\text{for all \(\chi\)}.
\]
Recall that already at all finite singularities \(\alpha\) in \(D\),
\(\mathcal{F}\) has
\[
  e_1(\alpha,\mathbf{1},\mathcal{F})\geq
  e_1(\alpha,\chi,\mathcal{F})\quad\text{for all \(\chi\)}.
\]

\sourcepara{5.2.4.3}
We will show this is impossible if \(\mathcal{F}\) in \(\mathcal{T}_\ell\) is
cohomologically rigid. First, we remark that on \(\mathbb{A}^1-D\),
\(\mathcal{F}\) is both irreducible and nontrivial (just because it lies in
\(\mathcal{T}_\ell\)). Therefore, denoting by
\[
  j:\mathbb{A}^1-D\to\mathbb{P}^1
\]
the inclusion, we have
\[
  \mathrm{H}^0(\mathbb{P}^1,j_*(\mathcal{F}\vert\mathbb{A}^1-D))=0
  =
  \mathrm{H}^2(\mathbb{P}^1,j_*(\mathcal{F}\vert\mathbb{A}^1-D)),
\]
and hence we have the inequality
\[
  \chi(\mathbb{P}^1,j_*(\mathcal{F}\vert\mathbb{A}^1-D))\leq 0.
\]

\sourcepara{5.2.4.4}
On the other hand, we are given that \(\mathcal{F}\) is cohomologically rigid,
i.e.,
\[
  \chi(\mathbb{P}^1,j_*\underline{\operatorname{End}}(\mathcal{F}\vert\mathbb{A}^1-D))=2.
\]

\sourcepara{5.2.4.5}
The idea now is to use the Euler--Poincaré formula to make explicit both of
these Euler characteristics. Consider any \(\mathcal{F}\) in
\(\mathcal{T}_\ell\) which is lisse on \(\mathbb{A}^1-D\). Because
\(\mathcal{F}\) is tame,
\[
\begin{aligned}
  &\chi(\mathbb{P}^1,j_*\underline{\operatorname{End}}(\mathcal{F}\vert\mathbb{A}^1-D))\\
  &=
  \chi(\mathbb{A}^1-D,\underline{\operatorname{End}}(\mathcal{F}))
  +\sum_{\alpha\in D\amalg\{\infty\}}
  \dim(\underline{\operatorname{End}}(\mathcal{F}))^{I(\alpha)}\\
  &=
  (1-\operatorname{Card}(D))r(\mathcal{F})^2
  +\sum_{\alpha\in D\amalg\{\infty\}}\sum_\chi\sum_i
  (e_i(\alpha,\chi,\mathcal{F}))^2,
\end{aligned}
\]
the last equality thanks to \hyperref[lem:3.1.15]{Lemma~3.1.15}.

\sourcepara{5.2.4.6}
Now at each \(\alpha\) in \(D\amalg\{\infty\}\), denote by
\(\chi_{\alpha,\mathcal{F}}\) a choice of character for which
\(e_1(\alpha,\chi,\mathcal{F})\) is maximal. Since the \(e_i\) are a
decreasing sequence, we have
\[
  e_1(\alpha,\chi_{\alpha,\mathcal{F}},\mathcal{F})
  \geq e_i(\alpha,\rho,\mathcal{F})\quad
  \text{for all \(\rho\) and for all \(i\)}.
\]
Therefore we have the inequality
\[
  \sum_\chi\sum_i(e_i(\alpha,\chi,\mathcal{F}))^2
  \leq
  \sum_\chi\sum_i
  e_1(\alpha,\chi_{\alpha,\mathcal{F}},\mathcal{F})
  e_i(\alpha,\chi,\mathcal{F}).
\]
But
\[
\begin{aligned}
  &\sum_\chi\sum_i
  e_1(\alpha,\chi_{\alpha,\mathcal{F}},\mathcal{F})
  e_i(\alpha,\chi,\mathcal{F})\\
  &=
  e_1(\alpha,\chi_{\alpha,\mathcal{F}},\mathcal{F})
  \sum_\chi\sum_i e_i(\alpha,\chi,\mathcal{F})\\
  &=e_1(\alpha,\chi_{\alpha,\mathcal{F}},\mathcal{F})r(\mathcal{F}),
\end{aligned}
\]
so we get
\[
\begin{aligned}
  &\chi(\mathbb{P}^1,j_*\underline{\operatorname{End}}(\mathcal{F}\vert\mathbb{A}^1-D))\\
  &\leq
  (1-\operatorname{Card}(D))r(\mathcal{F})^2
  +
  \sum_{\alpha\in D\amalg\{\infty\}}
  e_1(\alpha,\chi_{\alpha,\mathcal{F}},\mathcal{F})r(\mathcal{F}).
\end{aligned}
\]
This we record as

\sourceheading[ineq:5.2.4.7]{Basic Inequality 5.2.4.7}
For any \(\mathcal{F}\) in \(\mathcal{T}_\ell\) which is lisse on
\(\mathbb{A}^1-D\), we have
\[
\begin{aligned}
  &\chi(\mathbb{P}^1,j_*\underline{\operatorname{End}}(\mathcal{F}\vert\mathbb{A}^1-D))\\
  &\leq
  r(\mathcal{F})\bigl[
  (1-\operatorname{Card}(D))r(\mathcal{F})
  +
  \sum_{\alpha\in D\amalg\{\infty\}}
  e_1(\alpha,\chi_{\alpha,\mathcal{F}},\mathcal{F})
  \bigr].
\end{aligned}
\]

\sourcepara{5.2.4.8}
Now suppose that
\[
  (**)\quad
  \text{for every \(\alpha\) in \(D\amalg\{\infty\}\), we have }
  e_1(\alpha,\mathbf{1},\mathcal{F})\geq
  e_1(\alpha,\chi,\mathcal{F})\text{ for all \(\chi\)}.
\]
Then \(e_1(\alpha,\chi_{\alpha,\mathcal{F}},\mathcal{F})\) is just
\(e_1(\alpha,\mathbf{1},\mathcal{F})\), and we find
\[
\begin{aligned}
  &(1/r(\mathcal{F}))
  \chi(\mathbb{P}^1,j_*\underline{\operatorname{End}}(\mathcal{F}\vert\mathbb{A}^1-D))\\
  &\leq
  (1-\operatorname{Card}(D))r(\mathcal{F})
  +
  \sum_{\alpha\in D\amalg\{\infty\}}
  e_1(\alpha,\mathbf{1},\mathcal{F}).
\end{aligned}
\]
But the Euler--Poincaré formula for \(j_*\mathcal{F}\) gives
\[
  \chi(\mathbb{P}^1,j_*(\mathcal{F}\vert\mathbb{A}^1-D))
  =
  (1-\operatorname{Card}(D))r(\mathcal{F})
  +
  \sum_{\alpha\in D\amalg\{\infty\}}
  e_1(\alpha,\mathbf{1},\mathcal{F}).
\]
Thus for \(\mathcal{F}\) in \(\mathcal{T}_\ell\) satisfying \((**)\), we get
the inequality
\[
  \chi(\mathbb{P}^1,j_*\underline{\operatorname{End}}(\mathcal{F}\vert\mathbb{A}^1-D))
  \leq
  r(\mathcal{F})\chi(\mathbb{P}^1,j_*(\mathcal{F}\vert\mathbb{A}^1-D)).
\]
But \(\chi(\mathbb{P}^1,j_*(\mathcal{F}\vert\mathbb{A}^1-D))\leq 0\) for any
\(\mathcal{F}\) in \(\mathcal{T}_\ell\), so we get
\[
  \chi(\mathbb{P}^1,j_*\underline{\operatorname{End}}(\mathcal{F}\vert\mathbb{A}^1-D))
  \leq 0,
\]
which is impossible if \(\mathcal{F}\) is cohomologically rigid, i.e., if
\[
  \chi(\mathbb{P}^1,j_*\underline{\operatorname{End}}(\mathcal{F}\vert\mathbb{A}^1-D))
  =2.
\]

\sourcepara{5.2.4.9}
This shows that if we choose a character \(\chi\) so as to maximize the
dimension of the \(I(\infty)^{\mathrm{tame}}\) invariants in
\(\mathcal{L}_\chi\otimes\mathcal{F}\), then any such \(\chi\) is nontrivial,
provided \(\mathcal{F}\) is cohomologically rigid and satisfies the hypothesis
\((*)\) of \hyperref[thm:5.2.3]{Main Theorem~5.2.3}.

\sourcepara{5.2.4.10}
We must now show that for any such choice of \(\chi\), the sheaf
\(\mathcal{G}:=\operatorname{MC}_\chi(\mathcal{F})\) has strictly lower generic
rank: \(r(\mathcal{G})<r(\mathcal{F})\). For this, we use the rank formula
\hyperref[cor:3.3.7]{Corollary~3.3.7}, according to which, for any
\(\mathcal{F}\) in \(\mathcal{T}_\ell\),
\[
\begin{aligned}
  r(\mathcal{G})
  &=
  \sum_{\alpha\in D}
  \operatorname{rank}(\mathcal{F}(\alpha)/\mathcal{F}(\alpha)^{I(\alpha)})
  -
  \operatorname{rank}((\mathcal{F}(\infty)\otimes\mathcal{L}_\chi)^{I(\infty)})\\
  &=
  \operatorname{Card}(D)r(\mathcal{F})
  -
  \sum_{\alpha\in D}\operatorname{rank}(\mathcal{F}(\alpha)^{I(\alpha)})
  -
  \operatorname{rank}((\mathcal{F}(\infty)\otimes\mathcal{L}_\chi)^{I(\infty)})\\
  &=
  \operatorname{Card}(D)r(\mathcal{F})
  -
  \sum_{\alpha\in D}e_1(\alpha,\mathbf{1},\mathcal{F})
  -
  \operatorname{rank}((\mathcal{F}(\infty)\otimes\mathcal{L}_\chi)^{I(\infty)}).
\end{aligned}
\]

\sourcepara{5.2.4.11}
The Euler--Poincaré formula for \(\mathcal{F}\) says
\[
\begin{aligned}
  \chi(\mathbb{P}^1,j_*(\mathcal{F}\vert\mathbb{A}^1-D))
  &=
  (1-\operatorname{Card}(D))r(\mathcal{F})
  +
  \sum_{\alpha\in D\amalg\{\infty\}}e_1(\alpha,\mathbf{1},\mathcal{F})\\
  &=
  -\bigl[\operatorname{Card}(D)r(\mathcal{F})
  -\sum_{\alpha\in D}e_1(\alpha,\mathbf{1},\mathcal{F})\bigr]
  +r(\mathcal{F})+e_1(\infty,\mathbf{1},\mathcal{F}).
\end{aligned}
\]

\sourcepara{5.2.4.12}
Thus in our formula for \(r(\mathcal{G})\) we get, for any \(\mathcal{F}\) in
\(\mathcal{T}_\ell\),
\[
  r(\mathcal{G})=
  r(\mathcal{F})+e_1(\infty,\mathbf{1},\mathcal{F})
  -\chi(\mathbb{P}^1,j_*(\mathcal{F}\vert\mathbb{A}^1-D))
  -
  \operatorname{rank}((\mathcal{F}(\infty)\otimes\mathcal{L}_\chi)^{I(\infty)}).
\]

\sourcepara{5.2.4.13}
We rewrite this as
\[
\begin{aligned}
  r(\mathcal{F})-r(\mathcal{G})
  &=
  \chi(\mathbb{P}^1,j_*(\mathcal{F}\vert\mathbb{A}^1-D))
  +
  \operatorname{rank}((\mathcal{F}(\infty)\otimes\mathcal{L}_\chi)^{I(\infty)})
  -
  e_1(\infty,\mathbf{1},\mathcal{F}).
\end{aligned}
\]

\sourcepara{5.2.4.14}
So we need to show that, for our \(\mathcal{F}\), we have
\[
  \chi(\mathbb{P}^1,j_*(\mathcal{F}\vert\mathbb{A}^1-D))
  +
  \operatorname{rank}((\mathcal{F}(\infty)\otimes\mathcal{L}_\chi)^{I(\infty)})
  -
  e_1(\infty,\mathbf{1},\mathcal{F})>0.
\]

\sourcepara{5.2.4.15}
To show this, recall that for any \(\mathcal{F}\) in \(\mathcal{T}_\ell\) we
had the \hyperref[ineq:5.2.4.7]{Basic Inequality~5.2.4.7}
\[
\begin{aligned}
  &\chi(\mathbb{P}^1,j_*\underline{\operatorname{End}}(\mathcal{F}\vert\mathbb{A}^1-D))\\
  &\leq
  r(\mathcal{F})\bigl[
  (1-\operatorname{Card}(D))r(\mathcal{F})
  +
  \sum_{\alpha\in D\amalg\{\infty\}}
  e_1(\alpha,\chi_{\alpha,\mathcal{F}},\mathcal{F})
  \bigr].
\end{aligned}
\]
For \(\mathcal{F}\) cohomologically rigid and satisfying the hypothesis
\((*)\) of \hyperref[thm:5.2.3]{Main Theorem~5.2.3}, the left hand side is
\(2\) (\(\mathcal{F}\) is cohomologically rigid) and on the right each term
\(e_1(\alpha,\chi_{\alpha,\mathcal{F}},\mathcal{F})\) for \(\alpha\) in \(D\)
is equal to \(e_1(\alpha,\mathbf{1},\mathcal{F})\). Thus the Basic Inequality
gives
\[
\begin{aligned}
  2/r(\mathcal{F})
  &\leq
  (1-\operatorname{Card}(D))r(\mathcal{F})
  +e_1(\infty,\chi_{\infty,\mathcal{F}},\mathcal{F})
  +\sum_{\alpha\in D}e_1(\alpha,\mathbf{1},\mathcal{F})\\
  &=
  \chi(\mathbb{P}^1,j_*(\mathcal{F}\vert\mathbb{A}^1-D))
  +e_1(\infty,\chi_{\infty,\mathcal{F}},\mathcal{F})
  -e_1(\infty,\mathbf{1},\mathcal{F}).
\end{aligned}
\]
Since the left hand side (namely \(2/r(\mathcal{F})\)) is strictly positive, we
get
\[
  0<
  \chi(\mathbb{P}^1,j_*(\mathcal{F}\vert\mathbb{A}^1-D))
  +e_1(\infty,\chi_{\infty,\mathcal{F}},\mathcal{F})
  -e_1(\infty,\mathbf{1},\mathcal{F}).
\]
This is precisely the required inequality
\hyperref[par:5.2.4.14]{(5.2.4.14)}, for by the very choice of the character
\(\chi\) we have
\[
  e_1(\infty,\chi_{\infty,\mathcal{F}},\mathcal{F})
  =
  \operatorname{rank}((\mathcal{F}(\infty)\otimes\mathcal{L}_\chi)^{I(\infty)}).
\]

\section{Applications and interpretations of the main theorem}\label{sec:5.3}

\sourcepara{5.3.1}
We continue to work in the situation \hyperref[par:5.1.1]{(5.1.1)}. Fix a
dense open set \(\mathbb{A}^1-D\) in \(\mathbb{A}^1\), and a subgroup
\(\Gamma\) of the group of all continuous \(\overline{\mathbb{Q}}_\ell^\times\)-valued
characters of \(I(0)^{\mathrm{tame}}\). Let us denote by
\(\mathcal{T}_\ell(\mathbb{A}^1-D,\Gamma)\) the full subcategory of
\(\mathcal{T}_\ell\) consisting of the objects \(\mathcal{F}\) which are lisse
on \(\mathbb{A}^1-D\) and for which
\(\operatorname{LocMono}(\mathcal{F})\subset\Gamma\).

\sourcepara{5.3.2}
We can construct a graph whose vertices are the objects of
\(\mathcal{T}_\ell(\mathbb{A}^1-D,\Gamma)\), and in which there is an edge
joining two objects \(\mathcal{F}\) and \(\mathcal{G}\) if either of the
following conditions (a) or (b) holds:
\begin{enumerate}[label=(\alph*)]
\item \(\mathcal{F}\) and \(\mathcal{G}\) both have generic rank \(>1\), and
there is a lisse tame rank one \(\mathcal{L}\) on \(\mathbb{A}^1-D\), all of
whose local monodromies are in \(\Gamma\), such that
\(\mathcal{F}\cong\operatorname{MT}_{\mathcal{L}}(\mathcal{G})\), or
equivalently
\(\mathcal{G}\cong\operatorname{MT}_{\mathcal{L}^{-1}}(\mathcal{F})\).

\item there exists a nontrivial character \(\chi\) in \(\Gamma\) such that
\(\mathcal{F}\cong\operatorname{MC}_\chi(\mathcal{G})\), or equivalently
\(\mathcal{G}\cong\operatorname{MC}_{\chi^{-1}}(\mathcal{F})\).
\end{enumerate}
In terms of this graph, the Main Theorem says precisely

\sourceheading[thm:5.3.3]{Main Theorem 5.3.3 (= \hyperref[thm:5.2.1]{\textcolor{blue}{5.2.1 bis}})}
If \(\mathcal{F}\) in \(\mathcal{T}_\ell(\mathbb{A}^1-D,\Gamma)\) is
cohomologically rigid, then \(\mathcal{F}\) (as vertex of the graph
\hyperref[par:5.3.2]{(5.3.2)}) is connected to an object of rank one in
\(\mathcal{T}_\ell(\mathbb{A}^1-D,\Gamma)\), and its distance to such an
object is at most \(2(r(\mathcal{F})-1)\).

\section{Some open questions}\label{sec:5.4}

\sourcepara{5.4.1}
In this graph, any two objects which are connected have the same index of
rigidity
\(\chi(\mathbb{P}^1,j_*\underline{\operatorname{End}}(\mathcal{F}\vert\mathbb{A}^1-D))\),
by \hyperref[thm:4.3.10]{Theorem~4.3.10}. Let us say that an object
\(\mathcal{F}\) of \(\mathcal{T}_\ell(\mathbb{A}^1-D,\Gamma)\) is minimal if
for any object \(\mathcal{G}\) to which it is connected,
\(r(\mathcal{F})\leq r(\mathcal{G})\). Obviously every object is connected to a
minimal object which has the same index of rigidity. The
\hyperref[thm:5.3.3]{Main Theorem~5.3.3} states that the minimal
cohomologically rigid objects are exactly those of generic rank one.

\sourcepara{5.4.2}
But the index of rigidity, a priori even and \(\leq 2\), can be any even
integer \(\leq 2\). (For instance, the pullback by \(x\mapsto x^n\) of the rank
two local system on \(\mathbb{P}^1-\{0,1,\infty\}\) attached to the
differential equation for the Gauss hypergeometric function
\(F(1/2,1/2,1;x)\) has index of rigidity \(4-2n\).) What are the minimal
objects \(\mathcal{F}\) whose index of rigidity is some given even integer
\(\leq 0\)? Already the simplest case of this question, index of rigidity
\(=0\), is far from understood. For instance, is it true that if an object is
not minimal, then it is at distance at most two from an object of strictly
lower rank? Are there minimal objects with index of rigidity \(=0\) of
arbitrarily large rank?

\sourcepara{5.4.3}
Here is a ``candidate'' for a construction of minimal objects with index of
rigidity \(=0\) of arbitrarily large rank on \(\mathbb{A}^1-\{3\text{
points}\}\). Over \(\mathbb{C}\) take an elliptic curve
\[
  y^2=(x-e_1)(x-e_2)(x-e_3),
\]
an integer \(n\geq 1\), and a primitive \(n\)'th root of unity \(\zeta\).
Denote by \(\mathcal{F}_{n,\zeta}\) the (extension by zero to \(E\) of the)
rank \(n\) local system on \(E-\{0\}\) with local monodromy around \(0\) given
by the scalar \(\zeta\) constructed in the proof of
\hyperref[lem:1.4.4]{Lemma~1.4.4}. For a rank one local system \(\mathcal{L}\)
on \(E\) for which \(\mathcal{L}^{\otimes 2n}\) is nontrivial, there exists no
isomorphism (of local systems on \(E-\{0\}\))
\[
  \mathcal{L}\otimes\mathcal{F}_{n,\zeta}
  \cong[-1]^*(\mathcal{L}\otimes\mathcal{F}_{n,\zeta})
\]
(compare determinants). Let \(\pi:E\to\mathbb{P}^1\) be the map ``\(x\)''.
Then \(\mathcal{G}:=\pi_*(\mathcal{L}\otimes\mathcal{F}_{n,\zeta})\), for
\(\mathcal{L}\) as above, is an irreducible middle extension sheaf on
\(\mathbb{P}^1\) which on \(\mathbb{A}^1-\{e_1,e_2,e_3\}\) is a rank \(2n\)
local system with index of rigidity \(0\). Its local monodromies at the four
missing points are semisimple, with eigenvalues
\[
\begin{aligned}
  \text{at \(\infty\):}\quad
    &\text{each of the two square roots of \(\zeta\), repeated \(n\) times},\\
  \text{at \(e_i\):}\quad
    &\text{each of \(\pm 1\), repeated \(n\) times}.
\end{aligned}
\]
It is easy to show that any object of distance at most two from
\(\mathcal{G}\) has rank at least that of \(\mathcal{G}\). Is \(\mathcal{G}\)
minimal? If so, \(\mathcal{G}\) provides an example of a minimal object of rank
\(2n\), lisse on \(\mathbb{A}^1-\{e_1,e_2,e_3\}\), with index of rigidity
\(=0\).

\section[Universal families of rigids]{Existence of universal families of rigids with given local monodromy}
\label{sec:5.5}

\sourcepara{5.5.1}
Let \(k\) be an algebraically closed field, \(n\geq 2\) an integer,
\(\alpha_1,\alpha_2,\ldots,\alpha_n\) a set of \(n\) distinct points of
\(\mathbb{A}^1(k)\), \(\ell\) a prime number invertible in \(k\), \(N\geq 1\)
an integer which is invertible in \(k\), and \(\zeta\) a primitive \(N\)'th root
of unity in \(k\). Let \(\mathcal{F}\) be an object of \(\mathcal{T}_\ell\)
which is lisse on
\(\mathbb{A}^1-\{\alpha_1,\alpha_2,\ldots,\alpha_n\}\), cohomologically rigid,
and such that all eigenvalues of all local monodromies of \(\mathcal{F}\) are
\(N\)'th roots of unity.

\sourcepara{5.5.2}
Given the data \((n,N,\ell)\), we form the ring
\[
  R_{N,\ell}:=\mathbb{Z}[\zeta_N,1/N\ell],
  \quad \zeta_N:=\text{a primitive \(N\)'th root of unity}
\]
(i.e., \(R_{N,\ell}\) is
\(\mathbb{Z}[1/N\ell][X]/(\Phi_N(X))\), with \(\Phi_N(X)\) in
\(\mathbb{Z}[X]\) the \(N\)'th cyclotomic polynomial). We denote by
\(E=E_N\) the fraction field of \(R_{N,\ell}\): thus \(E\) is the cyclotomic
field \(\mathbb{Q}(\zeta_N)\). We fix an embedding
\[
  R_{N,\ell}\to\overline{\mathbb{Q}}_\ell,
\]
i.e., we fix a primitive \(N\)'th root of unity in
\(\overline{\mathbb{Q}}_\ell\). We denote by \(\lambda\) the induced place of
the ``abstract'' field \(E\), and by \(E_\lambda\) the \(\lambda\)-adic
completion of \(E\).

\sourcepara{5.5.3}
Denote by \(S_{N,n,\ell}\) the ring
\[
  S_{N,n,\ell}:=
  R_{N,\ell}[T_1,\ldots,T_n][1/\Delta],
  \quad
  \Delta:=\prod_{i\neq j}(T_i-T_j).
\]
There is a unique ring homomorphism
\[
  \varphi:S_{N,n,\ell}\to k
\]
for which \(\varphi(\zeta_N)=\zeta\) and for which
\(\varphi(T_i)=\alpha_i\) for \(1\leq i\leq n\).

Over \(S_{N,n,\ell}\), we have \(\mathbb{A}^1\), with its \(n\) disjoint
sections \(\{T_1,\ldots,T_n\}\). We denote by
\[
  j:(\mathbb{A}^1-\{T_1,\ldots,T_n\})_{S_{N,n,\ell}}
  \to
  (\mathbb{A}^1)_{S_{N,n,\ell}}
\]
the inclusion.

\sourceheading[thm:5.5.4]{Theorem 5.5.4}
In the situation \hyperref[par:5.5.1]{(5.5.1)}, we have
\begin{enumerate}[label=(\alph*)]
\item There exists a lisse \(E_\lambda\)-sheaf
\[
  \mathcal{F}\text{ on }
  (\mathbb{A}^1-\{T_1,\ldots,T_n\})_{S_{N,n,\ell}}
\]
which, after the base change \(\varphi:S_{N,n,\ell}\to k\) and the extension
of scalars \(E_\lambda\to\overline{\mathbb{Q}}_\ell\) becomes (the restriction
to \(\mathbb{A}^1-\{\alpha_1,\alpha_2,\ldots,\alpha_n\}\) of)
\(\mathcal{F}\).

\item (``mise pour memoire'', cf.
\hyperref[prop:4.2.3]{Proposition~4.2.3},
\hyperref[lem:4.3.8]{Lemma~4.3.8}, and
\hyperref[prop:4.3.9]{Proposition~4.3.9}) The object \(j_*\mathcal{F}\) on
\((\mathbb{A}^1)_{S_{N,n,\ell}}\) is of formation compatible with arbitrary
change of base on \(S_{N,n,\ell}\), and is adapted to the stratification
\[
  ((\mathbb{A}^1-\{T_1,\ldots,T_n\})_{S_{N,n,\ell}},
  \{T_1,\ldots,T_n\}_{S_{N,n,\ell}})
  \text{ of }(\mathbb{A}^1)_{S_{N,n,\ell}}.
\]
The restriction of \(j_*\mathcal{F}\) to every geometric fibre of
\((\mathbb{A}^1)_{S_{N,n,\ell}}\) is (after extension of scalars
\(E_\lambda\to\overline{\mathbb{Q}}_\ell\)) a cohomologically rigid object of
\(\mathcal{T}_\ell\), all of whose local monodromies have all their eigenvalues
\(N\)'th roots of unity.

\item The lisse \(E_\lambda\)-sheaf \(\mathcal{F}\) on
\((\mathbb{A}^1-\{T_1,\ldots,T_n\})_{S_{N,n,\ell}}\) is pure of some integer
weight \(w\), \(0\leq w\leq\operatorname{rank}(\mathcal{F})-1\), and all its
characteristic polynomials of Frobenius at all finite field-valued points have
coefficients in the subring \(\mathbb{Z}[\zeta_N]\) of \(E_\lambda\) (a subring
via the given embedding of \(R_{N,\ell}\) into
\(\overline{\mathbb{Q}}_\ell\)).

\item For any prime number \(\ell_1\), and any embedding
\[
  \lambda_1:\mathbb{Z}[\zeta_N]\to\overline{\mathbb{Q}}_{\ell_1},
\]
there exists on
\((\mathbb{A}^1-\{T_1,\ldots,T_n\})_{S_{N,n,\ell_1}}\) a lisse
\(E_{\lambda_1}\)-sheaf \(\mathcal{F}_{\lambda_1}\) which satisfies (2) above
with \(\ell\) replaced by \(\ell_1\), which is pure of the same weight \(w\),
which at every finite field-valued point has characteristic polynomial of
Frobenius with coefficients in the subring \(\mathbb{Z}[\zeta_N]\) of
\(E_{\lambda_1}\). Moreover, at any such finite field valued point where
\(\ell\) is invertible, the characteristic polynomial of
\(\mathcal{F}_{\lambda_1}\) is equal to that of \(\mathcal{F}\) (equality in
the common subring \(\mathbb{Z}[\zeta_N]\)).
\end{enumerate}

\sourceheading[proofhead:5.5.5]{5.5.5 Construction-proof of \hyperref[thm:5.5.4]{\textcolor{blue}{Theorem~5.5.4}}}

\sourcepara{5.5.5.1}
We proceed by induction on the generic rank \(r(\mathcal{F})\) of
\(\mathcal{F}\). Suppose first that \(r(\mathcal{F})=1\). Then on
\(\mathbb{A}^1-\{\alpha_1,\alpha_2,\ldots,\alpha_n\}\), \(\mathcal{F}\) is a
lisse sheaf of rank one, all of whose local monodromies have order dividing
\(N\). Denote by \(\chi_i\) the character of order dividing \(N\) of
\(I(\alpha_i)^{\mathrm{tame}}\) given by \(\mathcal{F}\) at \(\alpha_i\). Thus
\(\chi_i\) is a \(E_\lambda^\times\)-valued character of
\[
  I(\alpha_i)^{\mathrm{tame}}/NI(\alpha_i)^{\mathrm{tame}}\cong\mu_N(k),
\]
the isomorphism given by the action of \(\mu_N(k)\) as Galois group of the
finite étale Galois connected covering of
\(\mathbb{A}^1-\{\alpha_i\}\) of equation \(y^N=x-\alpha_i\), on which
\(\xi\) in \(\mu_N(k)\) acts by \((x,y)\mapsto(x,\xi y)\). Thus the translated
Kummer sheaf \(\mathcal{L}_{\chi_i(x-\alpha_i)}\) on
\(\mathbb{A}^1-\{\alpha_i\}\) is a lisse rank one \(E_\lambda\)-sheaf on
\(\mathbb{A}^1-\{\alpha_i\}\), and it has the same local monodromy at
\(\alpha_i\) that \(\mathcal{F}\) does.

\sourcepara{5.5.5.2}
Consider the lisse, rank one \(E_\lambda\)-sheaf
\[
  \bigotimes_i\mathcal{L}_{\chi_i(x-\alpha_i)}
  \text{ on }\mathbb{A}^1-\{\alpha_1,\alpha_2,\ldots,\alpha_n\}.
\]
It has the same local monodromy as \(\mathcal{F}\) at each point
\(\alpha_i\), and both it and \(\mathcal{F}\) are tame at \(\infty\).
Therefore
\[
  \mathcal{F}\cong
  \bigotimes_i\mathcal{L}_{\chi_i(x-\alpha_i)}
  \text{ on }\mathbb{A}^1-\{\alpha_1,\alpha_2,\ldots,\alpha_n\},
\]
because their ratio is lisse and tame on \(\mathbb{A}^1\), but
\(\pi_1^{\mathrm{tame}}(\mathbb{A}^1)=0\).

\sourcepara{5.5.5.3}
By means of the fixed choices of primitive \(N\)'th roots of unity \(\zeta_N\)
in \(R_{N,\ell}\subset S_{N,n,\ell}\) and \(\zeta\) in \(k\), we may identify
the groups \(\mu_N(R_{N,\ell})=\mu_N(S_{N,n,\ell})\) and \(\mu_N(k)\). This
allows us to view \(\chi_i\) as a character of the group
\(\mu_N(S_{N,n,\ell})\) which is the Galois group of the finite étale Galois
connected covering of \((\mathbb{A}^1-\{T_i\})_{S_{N,n,\ell}}\) of equation
\[
  y^N=x-T_i,
\]
on which \(\xi\) in \(\mu_N(S_{N,n,\ell})\) acts by
\((x,y)\mapsto(x,\xi y)\). Thus we may speak of the translated Kummer sheaf
\(\mathcal{L}_{\chi_i(x-T_i)}\) as a lisse \(E_\lambda\)-sheaf of rank one on
\((\mathbb{A}^1-\{T_i\})_{S_{N,n,\ell}}\).

\sourcepara{5.5.5.4}
We now define \(\mathcal{F}\) to be
\[
  \mathcal{F}:=
  \bigotimes_i\mathcal{L}_{\chi_i(x-T_i)}
  \text{ on }(\mathbb{A}^1-\{T_1,\ldots,T_n\})_{S_{N,n,\ell}}.
\]
It is obvious that (1) (and hence (2) also) and (3) are satisfied, with the
weight \(w=0\). Since the characters \(\chi_i\) have values in the subgroup
\(\mu_N(\mathbb{Z}[\zeta_N])\) of
\(\overline{\mathbb{Q}}_\ell^\times\), for any prime number \(\ell_1\), and
any embedding
\[
  \lambda_1:\mathbb{Z}[\zeta_N]\to\overline{\mathbb{Q}}_{\ell_1},
\]
we can view the \(\chi_i\) as \(E_{\lambda_1}^\times\)-valued characters, and
define \(\mathcal{F}_{\lambda_1}\) by the same recipe as above, viewing now
each \(\mathcal{L}_{\chi_i(x-T_i)}\) as a lisse rank one
\(E_{\lambda_1}\)-sheaf on
\((\mathbb{A}^1-\{T_i\})_{S_{N,n,\ell_1}}\). It is obvious that (4) is now
satisfied. (One should remark that, fibre by fibre, this
\(\mathcal{F}_{\lambda_1}\) does indeed have at least two finite singularities,
since along the section \(T_i\) its local monodromy is of exactly the same
order as was the local monodromy of the original \(\mathcal{F}\) at
\(\alpha_i\), namely the order of the character \(\chi_i\); and the assumption
that \(\mathcal{F}\) is in \(\mathcal{T}_\ell\) guarantees that \(\chi_i\) is
nontrivial for at least two distinct values of \(i\).) This concludes the
construction-proof in the case of generic rank one.

\sourcepara{5.5.5.5}
We now explain how to pass to the general case. Thus let \(\mathcal{F}\) in
\(\mathcal{T}_\ell\) be lisse on
\(\mathbb{A}^1-\{\alpha_1,\alpha_2,\ldots,\alpha_n\}\), cohomologically rigid,
with all local monodromy eigenvalues \(N\)'th roots of unity, and with generic
rank \(r(\mathcal{F})\geq 2\).

\sourcepara{5.5.5.6}
According to the main theorem \hyperref[thm:5.2.1]{Main Theorem~5.2.1}, there
exist
\begin{enumerate}[label=(\arabic*)]
\item an object \(\mathcal{G}\) in \(\mathcal{T}_\ell\), lisse on
\(\mathbb{A}^1-\{\alpha_1,\alpha_2,\ldots,\alpha_n\}\), with all local
monodromy eigenvalues \(N\)'th roots of unity, with
\(r(\mathcal{G})<r(\mathcal{F})\),

\item a lisse rank one \(\overline{\mathbb{Q}}_\ell\)-sheaf \(\mathcal{L}\) on
\(\mathbb{A}^1-\{\alpha_1,\alpha_2,\ldots,\alpha_n\}\), with all local
monodromies of order dividing \(N\),

\item a nontrivial \(\overline{\mathbb{Q}}_\ell\)-valued character \(\chi\) of
\(\pi_1(\mathbb{G}_m)^{\mathrm{tame}}\) of order dividing \(N\),

\item an isomorphism
\[
  \mathcal{F}\cong\operatorname{MT}_{\mathcal{L}}(\operatorname{MC}_\chi(\mathcal{G})).
\]
\end{enumerate}

\sourcepara{5.5.5.7}
By applying the argument given in the rank one case to \(\mathcal{L}\) on
\(\mathbb{A}^1-\{\alpha_1,\alpha_2,\ldots,\alpha_n\}\), and to
\(\mathcal{L}_\chi\) on \(\mathbb{G}_m=\mathbb{A}^1-\{0\}\), we construct
lisse rank one \(E_\lambda\)-sheaves \(\mathcal{L}\) on
\((\mathbb{A}^1-\{T_1,\ldots,T_n\})_{S_{N,n,\ell}}\) and
\(\mathcal{L}_\chi\) on \((\mathbb{A}^1-\{0\})_{S_{N,n,\ell}}\) which after
the base change \(\varphi:S_{N,n,\ell}\to k\) and the extension of scalars
\(E_\lambda\to\overline{\mathbb{Q}}_\ell\) become \(\mathcal{L}\) and
\(\mathcal{L}_\chi\) respectively. The conclusions (2), (3), and (4) of the
theorem hold for \(\mathcal{L}\) on
\((\mathbb{A}^1-\{T_1,\ldots,T_n\})_{S_{N,n,\ell}}\). These same conclusions
hold also for \(\mathcal{L}_\chi\) on
\((\mathbb{A}^1-\{0\})_{S_{N,n,\ell}}\) provided we delete the phrase ``in
\(\mathcal{T}_\ell\)'' from (2).

\sourcepara{5.5.5.8}
By induction, we may assume there exists a lisse \(E_\lambda\)-sheaf
\(\mathcal{G}\) on
\((\mathbb{A}^1-\{T_1,\ldots,T_n\})_{S_{N,n,\ell}}\) for which after the base
change \(\varphi:S_{N,n,\ell}\to k\) and the extension of scalars
\(E_\lambda\to\overline{\mathbb{Q}}_\ell\) becomes \(\mathcal{G}\), and for
which (2), (3), and (4) hold. We define \(\mathcal{F}\) on
\((\mathbb{A}^1-\{T_1,\ldots,T_n\})_{S_{N,n,\ell}}\) by
\[
  \mathcal{F}[1]:=
  \mathcal{L}\otimes
  j^*(j_*\mathcal{G}[1]*_{\mathrm{mid}+}j_{0*}\mathcal{L}_\chi[1]),
\]
where
\[
  j:(\mathbb{A}^1-\{T_1,\ldots,T_n\})_{S_{N,n,\ell}}\to
  (\mathbb{A}^1)_{S_{N,n,\ell}},
\]
and
\[
  j_0:(\mathbb{A}^1-\{0\})_{S_{N,n,\ell}}\to
  (\mathbb{A}^1)_{S_{N,n,\ell}}
\]
denote the inclusions.

\sourcepara{5.5.5.9}
That (1) holds results from \hyperref[par:4.3.2]{(4.3.2)},
\hyperref[par:4.3.3]{(4.3.3)}, and \hyperref[thm:4.3.10]{Theorem~4.3.10}. As
already remarked, (2) for \(\mathcal{F}\) is automatic once we have (1) (use
\hyperref[thm:4.3.10]{Theorem~4.3.10} to see that (after extension of scalars
\(E_\lambda\to\overline{\mathbb{Q}}_\ell\))
\(\mathcal{F}\otimes\mathcal{L}^{-1}\), and hence \(\mathcal{F}\), whose
generic rank is \(\geq 2\), is fibre-by-fibre in \(\mathcal{T}_\ell\)). To
prove (3) for \(\mathcal{F}\), it suffices to do so for
\((j_*\mathcal{G}[1]*_{\mathrm{mid}+}j_{0*}\mathcal{L}_\chi[1])[-1]\), because
we already know that (3) holds for \(\mathcal{L}\) itself, with weight \(w=0\).

\sourcepara{5.5.5.10}
We begin the proof of (3). Let \(\mathbb{F}\) be a finite field in which
\(N\ell\) is invertible, and \(\rho:S_{N,n,\ell}\to\mathbb{F}\) a ring
homomorphism. Pick a point \(\tau\) in
\(\mathbb{A}^1(\mathbb{F})-\{\rho(T_1),\ldots,\rho(T_n)\}\). Denote by \(U\)
the open set
\[
  U:=(\mathbb{A}^1-\{\rho(T_1),\ldots,\rho(T_n),\tau\})_{\mathbb{F}}
\]
of \(\mathbb{A}^1_{\mathbb{F}}\). According to
\hyperref[recipe:5.1.7]{Explicit Recipes~5.1.7}, the stalk of
\((j_*\mathcal{G}[1]*_{\mathrm{mid}+}j_{0*}\mathcal{L}_\chi[1])[-1]\) at the
geometric point \((\rho,\tau,\overline{\mathbb{F}})\) is the image of the
``forget supports'' map
\[
  \mathrm{H}^1_c(U\otimes\overline{\mathbb{F}},
  \mathcal{G}\otimes\mathcal{L}_{\chi(\tau-x)})
  \to
  \mathrm{H}^1(U\otimes\overline{\mathbb{F}},
  \mathcal{G}\otimes\mathcal{L}_{\chi(\tau-x)}).
\]
By induction, \(\mathcal{G}\) is pure weight \(w\),
\(0\leq w\leq\operatorname{rank}(\mathcal{G})-1\), while
\(\mathcal{L}_\chi\) is pure of weight zero. By
\cite[3.2.3 and remark following its statement]{De-WeilII} the image of the
``forget supports'' map is precisely the weight \(=w+1\) quotient of the group
\[
  \mathrm{H}^1_c(U\otimes\overline{\mathbb{F}},
  \mathcal{G}\otimes\mathcal{L}_{\chi(\tau-x)}),
\]
which is a priori of weight \(\leq w+1\). Now the cohomology groups
\[
  \mathrm{H}^i_c(U\otimes\overline{\mathbb{F}},
  \mathcal{G}\otimes\mathcal{L}_{\chi(\tau-x)})=0
  \quad\text{for \(i\neq 1\)}
\]
(for \(i=0\) because we are lisse on an open curve, for \(i=2\) because in
addition we are geometrically irreducible and nontrivial). Therefore the
characteristic polynomial of Frobenius \(\operatorname{Frob}_{\rho,\tau,\mathbb{F}}\)
on
\(\mathrm{H}^1_c(U\otimes\overline{\mathbb{F}},
\mathcal{G}\otimes\mathcal{L}_{\chi(\tau-x)})\) is, by the Lefschetz Trace
Formula \cite{Gro-FL}, equal to the \(L\)-function of the lisse sheaf
\(\mathcal{G}\otimes\mathcal{L}_{\chi(\tau-x)}\) on \(U\):
\[
\begin{aligned}
  L(U,\mathcal{G}\otimes\mathcal{L}_{\chi(\tau-x)},T)
  &=
  \det(1-T\operatorname{Frob}_{\rho,\tau,\mathbb{F}}\mid
  \mathrm{H}^1_c(U\otimes\overline{\mathbb{F}},
  \mathcal{G}\otimes\mathcal{L}_{\chi(\tau-x)})).
\end{aligned}
\]
Therefore the characteristic polynomial of Frobenius
\(\operatorname{Frob}_{\rho,\tau,\mathbb{F}}\) on
\((j_*\mathcal{G}[1]*_{\mathrm{mid}+}j_{0*}\mathcal{L}_\chi[1])[-1]\) is the
``pure of weight \(w+1\) part'' of this \(L\)-function.

\sourcepara{5.5.5.11}
We now consider more closely this \(L\)-function. Because both \(\mathcal{G}\)
and \(\mathcal{L}_\chi\) have all their characteristic polynomials of Frobenius
with coefficients in the subring \(\mathbb{Z}[\zeta_N]\) of
\(\overline{\mathbb{Q}}_\ell\), the Euler product for this \(L\)-function shows
that the \(L\)-function itself lies in \(1+T(\mathbb{Z}[\zeta_N][[T]])\),
hence (being a polynomial) in \(1+T(\mathbb{Z}[\zeta_N][T])\), and that it is
entirely determined by all the individual characteristic polynomials of
Frobenius of both \(\mathcal{G}\) and \(\mathcal{L}_\chi\). If we factor
\(L(T)\) as
\[
  L(T)=\prod_j(1-\beta_jT),
\]
the reciprocal roots \(\beta_j\) are algebraic integers, which (as a set with
multiplicity) are stable under \(\operatorname{Aut}(\mathbb{C}/\mathbb{Q}(\zeta_N))\).
Taking the ``part of weight \(w+1\)'' of \(L\) produces the intrinsic divisor
of \(L(T)\) in which we keep precisely those \(\beta_j\) which, together with
all their conjugates, are algebraic integers, have complex absolute value
\(\sqrt{(\operatorname{Card}(\mathbb{F}))^{w+1}}\). This intrinsic divisor is
also stable under \(\operatorname{Aut}(\mathbb{C}/\mathbb{Q}(\zeta_N))\), so
has coefficients in
\(\mathbb{Q}(\zeta_N)\cap\{\text{algebraic integers}\}=\mathbb{Z}[\zeta_N]\).
Moreover, being intrinsic, it too is entirely determined by all the individual
characteristic polynomials of Frobenius of both \(\mathcal{G}\) and
\(\mathcal{L}_\chi\).

\sourcepara{5.5.5.12}
It remains to prove (4). Again by induction, all the objects
\(\mathcal{G}\), \(\mathcal{L}\), and \(\mathcal{L}_\chi\) have
\(\lambda_1\)-adic versions, say \(\mathcal{G}_{\lambda_1}\),
\(\mathcal{L}_{\lambda_1}\), and \(\mathcal{L}_{\chi_{\lambda_1}}\). Using
these, we define \(\mathcal{F}_{\lambda_1}\) out of them in precisely the same
way we defined \(\mathcal{F}\):
\[
  \mathcal{F}_{\lambda_1}[1]:=
  \mathcal{L}_{\lambda_1}\otimes
  j^*(j_*\mathcal{G}_{\lambda_1}[1]*_{\mathrm{mid}+}
  j_{0*}\mathcal{L}_{\chi_{\lambda_1}}[1]).
\]
Now repeat for
\[
  j^*(j_*\mathcal{G}_{\lambda_1}[1]*_{\mathrm{mid}+}
  j_{0*}\mathcal{L}_{\chi_{\lambda_1}}[1])[-1]
\]
the above discussion of characteristic polynomials of Frobenius. It will
establish that its characteristic polynomials of Frobenius have coefficients
in \(\mathbb{Z}[\zeta_N]\), and are entirely determined by all the individual
characteristic polynomials of Frobenius of both \(\mathcal{G}_{\lambda_1}\) and
\(\mathcal{L}_{\chi_{\lambda_1}}\) by exactly the same ``part of weight
\(w+1\) of an \(L\)-function'' rules as the characteristic polynomials of
\(\mathcal{F}\) were determined by those of \(\mathcal{G}\) and
\(\mathcal{L}_\chi\). This proves inductively that
\[
  j^*(j_*\mathcal{G}_{\lambda_1}[1]*_{\mathrm{mid}+}
  j_{0*}\mathcal{L}_{\chi_{\lambda_1}}[1])[-1]
\]
is pure of weight \(w+1\), has characteristic polynomials of Frobenius with
coefficients in \(\mathbb{Z}[\zeta_N]\), and that for any finite field-valued
point of \((\mathbb{A}^1-\{T_1,\ldots,T_n\})_{S_{N,n,\ell}}\) where
\(\ell_1\) is invertible,
\[
  j^*(j_*\mathcal{G}[1]*_{\mathrm{mid}+}j_{0*}\mathcal{L}_\chi[1])[-1]
  \quad\text{and}\quad
  j^*(j_*\mathcal{G}_{\lambda_1}[1]*_{\mathrm{mid}+}
  j_{0*}\mathcal{L}_{\chi_{\lambda_1}}[1])[-1]
\]
have the ``same'' (comparison in the common subring \(\mathbb{Z}[\zeta_N]\) of
\(\overline{\mathbb{Q}}_\ell\) and \(\overline{\mathbb{Q}}_{\ell_1}\))
characteristic polynomial of Frobenius. Because we already know that (4) holds
for \(\mathcal{L}\) and its \(\lambda_1\)-partner
\(\mathcal{L}_{\lambda_1}\), we may tensor by these to deduce that
\(\mathcal{F}\) and \(\mathcal{F}_{\lambda_1}\) also have the same
characteristic polynomials of Frobenius at any finite field-valued point of
\((\mathbb{A}^1-\{T_1,\ldots,T_n\})_{S_{N,n,\ell}}\) where \(\ell_1\) is
invertible.

\sourcepara{5.5.5.13}
It remains only to show that \(j_*\mathcal{F}_{\lambda_1}\) satisfies (2).
Using \hyperref[lem:4.3.8]{Lemma~4.3.8}, we get most of (2): it remains only
to show that \(j_*\mathcal{F}_{\lambda_1}\) is fibre-by-fibre in
\(\mathcal{T}_\ell\). But as \(\operatorname{rank}(\mathcal{F}_{\lambda_1})\geq 2\),
it is equivalent to show that
\[
  (j_*\mathcal{G}_{\lambda_1}[1]*_{\mathrm{mid}+}
  j_{0*}\mathcal{L}_{\chi_{\lambda_1}}[1])[-1]
\]
is fibre-by-fibre in \(\mathcal{T}_\ell\). This results from
\hyperref[thm:4.3.10]{Theorem~4.3.10} and
\hyperref[thm:4.3.11]{Theorem~4.3.11}.

\sourcepara{5.5.6}
We now discuss the weight \(w\) of \(\mathcal{F}\) constructed in
\hyperref[thm:5.5.4]{Theorem~5.5.4}. If \(\mathcal{F}\) has rank one, then
\(w=0\). When \(\mathcal{F}\) has rank \(r(\mathcal{F})\geq 2\), we can, by
\hyperref[thm:5.3.3]{Main Theorem~5.3.3}, connect \(\mathcal{F}\) to a rank
one object \(\mathcal{L}_0\) through a finite number of steps, each of which is
either an \(\operatorname{MT}_{\mathcal{L}}\) or an
\(\operatorname{MC}_\chi\), say
\[
  \mathcal{F}=
  \operatorname{MT}_{\mathcal{L}_d}\circ\operatorname{MC}_{\chi_d}
  \circ\cdots\circ
  \operatorname{MT}_{\mathcal{L}_1}\circ\operatorname{MC}_{\chi_1}
  (\mathcal{L}_0),
\]
in such a way that
\begin{enumerate}[label=(\arabic*)]
\item each application of an \(\operatorname{MC}_\chi\) strictly increases the
rank,
\item we have \((\chi_i)^N=1\) for \(i=1,\ldots,d\), but each
\(\chi_i\neq\mathbf{1}\),
\item we have \((\mathcal{L}_i)^{\otimes N}\cong\overline{\mathbb{Q}}_\ell\),
for \(i=0,\ldots,d\), but \(\mathcal{L}_i\) is non-constant for
\(i=0,\ldots,d-1\).
\end{enumerate}
There may be more than one such expression for \(\mathcal{F}\), but each
choice of such an expression gives rise to an \(\mathcal{F}\). Concretely, we
``thicken'' each \(\mathcal{L}_i\) to \(\mathcal{L}_i\), each
\(\mathcal{L}_{\chi_i}\) to \(\mathcal{L}_{\chi_i}\), and inductively define
\(\mathcal{F}_0:=\mathcal{L}_0\), \(\mathcal{F}_1,\ldots,\mathcal{F}_d:=\mathcal{F}\)
by (in the notations of \hyperref[par:5.5.5.8]{(5.5.5.8)})
\[
  \mathcal{F}_{i+1}[1]:=
  \mathcal{L}_i\otimes
  j^*(j_*\mathcal{F}_{i-1}[1]*_{\mathrm{mid}+}j_{0*}\mathcal{L}_{\chi_i}[1]),
\]
for \(i=0,\ldots,d-1\). The weight of the \(\mathcal{F}\) constructed this way
is the number \(d\) of steps of type \(\operatorname{MC}_\chi\).

\sourcepara{5.5.7}
We now examine the unicity of an \(\mathcal{F}\) given by
\hyperref[thm:5.5.4]{Theorem~5.5.4}. Denote by
\[
  \pi:(\mathbb{A}^1-\{T_1,\ldots,T_n\})_{S_{N,n,\ell}}\to S_{N,n,\ell}
\]
the projection. For any lisse \(E_\lambda\)-sheaf \(\mathcal{L}\) on
\(S_{N,n,\ell}\), \(\mathcal{F}\otimes\pi^*(\mathcal{L})\) works just as well
as \(\mathcal{F}\) in \hyperref[thm:5.5.4]{Theorem~5.5.4(1)}. This is the only
ambiguity in \(\mathcal{F}\).

\sourceheading[lem:5.5.7.1]{Lemma 5.5.7.1}
Hypotheses and notations as in \hyperref[thm:5.5.4]{Theorem~5.5.4}, suppose
\(\mathcal{F}\) and \(\mathcal{F}_1\) both satisfy
\hyperref[thm:5.5.4]{Theorem~5.5.4(1)}. Then there exists a lisse
\(E_\lambda\)-sheaf \(\mathcal{L}\) on \(S_{N,n,\ell}\), and an isomorphism
\[
  \mathcal{F}_1\cong\mathcal{F}\otimes\pi^*(\mathcal{L}).
\]

\begin{proof}
This follows from the Rigidity Corollary \hyperref[cor:5.5.7.3]{Corollary~5.5.7.3}
below, applied to \(S:=\operatorname{Spec}(S_{N,n,\ell})\), and
\(U:=(\mathbb{A}^1-\{T_1,\ldots,T_n\})_{S_{N,n,\ell}}\).
\end{proof}

\sourceheading[lem:5.5.7.2]{Lemma 5.5.7.2 (mise pour memoire)}
Let \(\ell\) be a prime number, \(S\) a normal connected noetherian
\(\mathbb{Z}[1/\ell]\)-scheme whose generic point has
characteristic zero, \(X/S\) a proper smooth curve with geometrically
connected fibres, \(D\subset X\) a closed subscheme which is finite étale over
S of degree \(d\geq 0\), \(U:=X-D\), and \(\pi:U\to S\) the structural map.
For any finite extension \(E_\lambda\) of \(\mathbb{Q}_\ell\), and any lisse
\(E_\lambda\)-sheaf \(\mathcal{F}\) on \(U\), the sheaves
\(\mathrm{R}^i\pi_!\mathcal{F}\) and \(\mathrm{R}^i\pi_*\mathcal{F}\) on \(S\)
are lisse, and of formation compatible with arbitrary change of base on \(S\).

\begin{proof}
Because \(S\) has generic characteristic zero, \(\mathcal{F}\) is automatically
tame on the geometric generic fibre, and the asserted result for
\(\mathrm{R}^i\pi_!\mathcal{F}\) is well-known, cf. \cite[4.7.1]{Ka-SE}. This
result, applied to \(\mathcal{F}^{\vee}\), yields the result for
\(\mathrm{R}^i\pi_*\mathcal{F}\) by duality.
\end{proof}

\sourceheading[cor:5.5.7.3]{Rigidity Corollary 5.5.7.3}
In the situation \hyperref[lem:5.5.7.2]{Lemma~5.5.7.2}, let \(\mathcal{F}\)
and \(\mathcal{G}\) be lisse \(E_\lambda\)-sheaves (respectively
\(\overline{\mathbb{Q}}_\ell\)-sheaves) on \(U\). The following conditions are
equivalent.
\begin{enumerate}[label=(\arabic*)]
\item For some geometric point \(s\) of \(S\), both \(\mathcal{F}\vert U_s\)
and \(\mathcal{G}\vert U_s\) are absolutely irreducible, and there exists an
isomorphism
\(\mathcal{F}\vert U_s\cong\mathcal{G}\vert U_s\) of lisse sheaves on \(U_s\).

\item For every geometric point \(s\) of \(S\), both
\(\mathcal{F}\vert U_s\) and \(\mathcal{G}\vert U_s\) are absolutely
irreducible, and there exists an isomorphism
\(\mathcal{F}\vert U_s\cong\mathcal{G}\vert U_s\) of lisse sheaves on \(U_s\).

\item Condition (2) holds, and there exists a lisse, rank one
\(E_\lambda\)-sheaf (respectively \(\overline{\mathbb{Q}}_\ell\)-sheaf)
\(\mathcal{L}\) on \(S\), and an isomorphism
\(\mathcal{G}\cong\mathcal{F}\otimes\pi^*(\mathcal{L})\) of lisse sheaves on
\(U\).
\end{enumerate}

\begin{proof}
It is trivial that (3) \(\Rightarrow\) (2) \(\Rightarrow\) (1). We now show
that (1) \(\Rightarrow\) (3). Consider the lisse sheaf
\(\mathcal{H}:=\underline{\operatorname{Hom}}(\mathcal{F},\mathcal{G})\) on
\(U\). By \hyperref[lem:5.5.7.2]{Lemma~5.5.7.2} applied to \(\mathcal{H}\),
\(\pi_*\underline{\operatorname{Hom}}(\mathcal{F},\mathcal{G})\) is lisse on
\(S\), of formation compatible with arbitrary change of base on \(S\). Taking
the fibre at \(s\), we see from (1) that
\(\mathcal{L}:=\pi_*\underline{\operatorname{Hom}}(\mathcal{F},\mathcal{G})\)
is lisse of rank one. We have a canonical map of lisse sheaves on \(U\),
\[
  \mathcal{F}\otimes\pi^*(\mathcal{L})\to\mathcal{G},
\]
which is an isomorphism on \(U_s\). But the kernel and cokernel of this map are
lisse sheaves on \(U\), so both must vanish. Therefore
\(\mathcal{F}\otimes\pi^*(\mathcal{L})\cong\mathcal{G}\). For every geometric
point \(s\) of \(S\), we get
\(\mathcal{F}\vert U_s\cong\mathcal{G}\vert U_s\) by passing to fibres. That
both \(\mathcal{F}\vert U_s\) and \(\mathcal{G}\vert U_s\) are absolutely
irreducible for every geometric point \(s\) results from
\hyperref[cor:4.2.6]{Corollary~4.2.6}.
\end{proof}

\sourcepara{5.5.8}
We now clarify the sense in which the \(\mathcal{F}\) constructed in
\hyperref[thm:5.5.4]{Theorem~5.5.4} is a universal family of rigids with the
same local monodromy as \(\mathcal{F}\). We must first specify the precise
meaning of ``same local monodromy as \(\mathcal{F}\)''. Thus let \(k_1\) be a
second algebraically closed field in which \(\ell N\) is invertible,
\(\zeta_1\) in \(k_1\) a primitive \(N\)'th root of unity, and
\(\beta_1,\ldots,\beta_n\) a set of \(n\geq 2\) distinct points of
\(\mathbb{A}^1(k_1)\). Let \(\mathcal{F}_1\) be a lisse
\(\overline{\mathbb{Q}}_\ell\)-sheaf on
\(\mathbb{A}^1_{k_1}-\{\beta_1,\ldots,\beta_n\}\) which is everywhere tame and
such that all eigenvalues of all local monodromies of \(\mathcal{F}_1\) are
\(N\)'th roots of unity. At any point
\(\beta\) in \(\{\beta_1,\ldots,\beta_n,\infty\}\), we can describe the local
monodromy of \(\mathcal{F}_1\) at \(\beta\) as follows. The tame inertia group
\(I(\beta)^{\mathrm{tame}}\) is canonically the pro-cyclic group
\[
  \prod_{\ell\neq\operatorname{char}(k)}\mathbb{Z}_\ell(1)(k_1)
  \cong\varprojlim_M\mu_M(k_1),
\]
the inverse limit taken over \(M\)'s invertible in \(k_1\). This group maps
onto \(\mu_N(k_1)\). By hypothesis, \(\mathcal{F}_1(\beta)\) as
\(I(\beta)^{\mathrm{tame}}\)-representation is canonically the direct sum of
representations of the form
\[
  (\text{a character of \(\mu_N(k_1)\)})\otimes
  (\text{a unipotent rep'n. of \(I(\beta)^{\mathrm{tame}}\)}).
\]
A unipotent representation of \(I(\beta)^{\mathrm{tame}}\) factors through a
unipotent representation of \(\mathbb{Z}_\ell(1)(k_1)\), and its isomorphism
class is the Jordan normal form of the action of any topological generator of
\(\mathbb{Z}_\ell(1)(k)\).

Therefore \(\mathcal{F}_1(\beta)\) as \(I(\beta)^{\mathrm{tame}}\)-representation
factors through the quotient
\[
  \varprojlim_\nu\mu_{\ell^\nu N}(k_1)
\]
of \(I(\beta)^{\mathrm{tame}}\). Moreover, for any choice
\(\gamma(\beta,\zeta_1)^{\mathrm{tame}}\) of a generator of this last group
which maps onto \(\zeta_1\) in \(\mu_N(k_1)\), the Jordan normal form (i.e.,
conjugacy class in \(\operatorname{GL}(\operatorname{rank}(\mathcal{F}_1),
\overline{\mathbb{Q}}_\ell)\)) of its action on \(\mathcal{F}_1(\beta)\) is
independent of the choice. By ``the local monodromy of \(\mathcal{F}_1\) at
\(\beta\), relative to \(\zeta_1\)'', we mean the Jordan normal form of the
action on \(\mathcal{F}_1(\beta)\) of any
\(\gamma(\beta,\zeta_1)^{\mathrm{tame}}\).

\sourceheading[prop:5.5.8.1]{Proposition 5.5.8.1}
Consider the situation \hyperref[par:5.5.1]{(5.5.1)}. Let \(k_1\) be a second
algebraically closed field in which \(\ell N\) is invertible, \(\zeta_1\) in
\(k_1\) a primitive \(N\)'th root of unity, and
\(\beta_1,\ldots,\beta_n\) a set of \(n\geq 2\) distinct
\(\mathbb{A}^1(k_1)\). Let \(\mathcal{F}_1\) be an object of
\(\mathcal{T}_\ell\) on \(\mathbb{A}^1_{k_1}\) which is lisse on
\(\mathbb{A}^1_{k_1}-\{\beta_1,\ldots,\beta_n\}\), cohomologically rigid, and
such that all eigenvalues of all local monodromies of \(\mathcal{F}_1\) are
\(N\)'th roots of unity. Suppose that
\begin{enumerate}[label=(\arabic*)]
\item for \(i=1,\ldots,n\), the local monodromy of \(\mathcal{F}_1\) at
\(\beta_i\), relative to \(\zeta_1\), is isomorphic to the local monodromy of
\(\mathcal{F}\) at \(\alpha_i\), relative to \(\zeta\),

\item the local monodromy of \(\mathcal{F}_1\) at \(\infty\), relative to
\(\zeta_1\), is isomorphic to the local monodromy of \(\mathcal{F}\) at
\(\infty\), relative to \(\zeta\).
\end{enumerate}

Denote by \(\mathcal{F}\) and \(\mathcal{F}_1\) on
\((\mathbb{A}^1-\{T_1,\ldots,T_n\})_{S_{N,n,\ell}}\) any choices of the lisse
\(E_\lambda\)-sheaves given by \hyperref[thm:5.5.4]{Theorem~5.5.4}, applied
to \(\mathcal{F}\) and to \(\mathcal{F}_1\) respectively. Denote by
\[
  \pi:(\mathbb{A}^1-\{T_1,\ldots,T_n\})_{S_{N,n,\ell}}\to S_{N,n,\ell}
\]
the projection. Then
\begin{enumerate}[label=(\arabic*)]
\item At any geometric point \(\varphi:S_{N,n,\ell}\to k_2\) of
\(S_{N,n,\ell}\), the restriction \(\mathcal{F}_{\varphi}\) of
\(\mathcal{F}\) to the fibre over \(\varphi\) is a cohomologically rigid object
of \(\mathcal{T}_\ell\), whose local monodromy at the point \(\varphi(T_i)\),
\(i=1,\ldots,n\), (respectively at \(\infty\)) relative to
\(\varphi(\zeta_N)\), is isomorphic to the local monodromy of \(\mathcal{F}\)
at \(\alpha_i\), \(i=1,\ldots,n\) (respectively at \(\infty\)), relative to
\(\zeta\).

\item Denote by \(s_1\) the geometric point
\(\varphi_1:S_{N,n,\ell}\to k_1\) with \(\varphi_1(\zeta_N)=\zeta_1\) and
\(\varphi_1(T_i)=\beta_i\) for \(i=1,\ldots,n\). The restriction
\(\mathcal{F}_{s_1}\) of \(\mathcal{F}\) to the fibre over \(s_1\) is
isomorphic to \(\mathcal{F}_1\).

\item There exists a lisse \(E_\lambda\)-sheaf \(\mathcal{L}\) on
\(S_{N,n,\ell}\) and an isomorphism
\[
  \mathcal{F}_1\cong\mathcal{F}\otimes\pi^*(\mathcal{L}).
\]
\end{enumerate}

\begin{proof}
Statement (1) results from \hyperref[thm:5.5.4]{Theorem~5.5.4(2)}, together
with \cite[1.7.8]{De-WeilII} (cf. \cite[4.7.2]{Ka-SE} and the proof of
\hyperref[thm:4.3.11]{Theorem~4.3.11}), which tells us that the local
monodromy of \(\mathcal{F}\) along each section \(T_i\) and along the section
\(\infty\) is ``the same'' on all geometric fibres. (Strictly speaking, we
should first extend scalars in the universal situation from
\(\mathbb{Z}[\zeta_N,1/N\ell]\) to
\(\mathbb{Z}[\text{all }\zeta_{\ell^\nu N},1/N\ell]\), choose a topological
generator of
\(\varprojlim_\nu\mu_{\ell^\nu N}(\mathbb{Z}[\text{all }\zeta_{\ell^\nu N},1/N\ell])\)
which maps to \(\zeta_N\), and use the fact that fibre by fibre this generator
gives a choice of \(\gamma(\beta,\zeta_1)^{\mathrm{tame}}\).) By (1), the
restriction \(\mathcal{F}_{s_1}\) of \(\mathcal{F}\) to the fibre over \(s_1\)
has the same local monodromy as \(\mathcal{F}_1\), at each of the points
\(\{\beta_1,\ldots,\beta_n,\infty\}\). Since \(\mathcal{F}_1\) is
cohomologically rigid, \(\mathcal{F}_1\) is isomorphic to
\(\mathcal{F}_{s_1}\). To get (3), apply the rigidity corollary
\hyperref[cor:5.5.7.3]{Corollary~5.5.7.3} to \(\mathcal{F}\) and to
\(\mathcal{F}_1\), and the geometric point \(s_1\).
\end{proof}

\section{Remark on braid groups}\label{sec:5.6}

The space \((\mathbb{A}^1-\{T_1,\ldots,T_n\})_{S_{N,n,\ell}}\) is the spec of
\[
  \mathbb{Z}[\zeta_N,1/N\ell][X,T_1,\ldots,T_n]
  [1/((\prod_{i\neq j}(T_i-T_j))(\prod_j(X-T_j)))].
\]
So if name the variables \(X\) as \(T_{n+1}\), we can think of this as the ring
\[
  \mathbb{Z}[\zeta_N,1/N\ell][T_1,\ldots,T_{n+1}]
  [1/(\prod_{i\neq j}(T_i-T_j))].
\]
Thought of this way, we see its fundamental group as an arithmetic-geometrical
version of the Artin braid group on \(n+1\) letters. And so we can restate the
first assertion of the theorem by the catch-phrase ``an irreducible local
system on \(\mathbb{A}^1\) minus \(n\) points which is cohomologically rigid
extends to a representation of the braid group on \(n+1\) letters''.

\section{Universal families without quasiunipotence}\label{sec:5.7}

\sourcepara{5.7.1}
We next explore the situation if we drop the hypothesis of quasiunipotence.
Thus, we let \(k\) be an algebraically closed field, \(n\geq 2\) an integer,
\(\alpha_1,\alpha_2,\ldots,\alpha_n\) a set of \(n\) distinct points of
\(\mathbb{A}^1(k)\), and \(\ell\) a prime number invertible in \(k\). Let
\(\mathcal{F}\) be a cohomologically rigid object of \(\mathcal{T}_\ell\) which
is lisse on \(\mathbb{A}^1-\{\alpha_1,\alpha_2,\ldots,\alpha_n\}\).

\sourcepara{5.7.2}
We denote by \(p\) the characteristic of \(k\). If \(p>0\), we denote by
\(\overline{\mathbb{F}}_p\) the algebraic closure of \(\mathbb{F}_p\) in \(k\),
and we denote by \(R\) the ring \(W(\overline{\mathbb{F}}_p)\) of Witt vectors
over \(\overline{\mathbb{F}}_p\). If \(\operatorname{char}(k)=0\), we denote
by \(R\) the subfield \(\mathbb{Q}(\text{all roots of unity})\) of \(k\). There
is a canonical ring homomorphism from \(R\) to \(k\), which for \(p=0\) is the
inclusion, and which for \(p>0\) is the composite
\[
  R\to R/pR=\overline{\mathbb{F}}_p\to k
\]
of reduction mod \(p\) and of the inclusion.

\sourcepara{5.7.3}
We denote by \(S_n\) the ring
\[
  S_n:=R[T_1,\ldots,T_n][1/\Delta],
  \quad
  \Delta:=\prod_{i\neq j}(T_i-T_j).
\]
There is a unique ring homomorphism
\[
  \varphi:S_n\to k
\]
which induces the canonical map on \(R\), and for which
\(\varphi(T_i)=\alpha_i\) for \(1\leq i\leq n\). Over \(S_n\), we have
\(\mathbb{A}^1\), with its \(n\) disjoint sections
\(\{T_1,\ldots,T_n\}\). We denote by
\[
  j:(\mathbb{A}^1-\{T_1,\ldots,T_n\})_{S_n}\to(\mathbb{A}^1)_{S_n}
\]
the inclusion.

A straightforward modification of the proof of the previous
\hyperref[thm:5.5.4]{Theorem~5.5.4} yields:

\sourceheading[thm:5.7.4]{Theorem 5.7.4}
In the situation \hyperref[par:5.7.1]{(5.7.1)}, we have
\begin{enumerate}[label=(\arabic*)]
\item There exists a lisse \(\overline{\mathbb{Q}}_\ell\)-sheaf
\[
  \mathcal{F}\text{ on }(\mathbb{A}^1-\{T_1,\ldots,T_n\})_{S_n}
\]
which, after the base change \(\varphi:S_n\to k\), becomes (the restriction to
\(\mathbb{A}^1-\{\alpha_1,\alpha_2,\ldots,\alpha_n\}\) of) \(\mathcal{F}\).

\item (``mise pour memoire'', cf.
\hyperref[prop:4.2.3]{Proposition~4.2.3},
\hyperref[lem:4.3.8]{Lemma~4.3.8},
\hyperref[prop:4.3.9]{Proposition~4.3.9}) The object \(j_*\mathcal{F}\) on
\((\mathbb{A}^1)_{S_n}\) is of formation compatible with arbitrary change of
base on \(S_{N,n,\ell}\), and is adapted to the stratification
\[
  ((\mathbb{A}^1-\{T_1,\ldots,T_n\})_{S_n},
  \{T_1,\ldots,T_n\}_{S_n})
  \text{ of }(\mathbb{A}^1)_{S_n}.
\]
At any geometric point \(\varphi:S_n\to k_2\) of \(S_n\), the restriction of
\(j_*\mathcal{F}\) to the fibre over \(\varphi\) is a cohomologically rigid
object of \(\mathcal{T}_\ell\), whose local monodromy at each point
\(\varphi(T_i)\) (respectively at \(\infty\)) is isomorphic to the local
monodromy of \(\mathcal{F}\) at \(\alpha_i\) (respectively at \(\infty\)).
\end{enumerate}

\sourceheading[rem:5.7.5]{Remark 5.7.5}
The space \((\mathbb{A}^1-\{T_1,\ldots,T_n\})_{S_n}\) is the spec of
\[
  R[T_1,\ldots,T_{n+1}][1/(\prod_{i\neq j}(T_i-T_j))].
\]
Its fundamental group is a (less arithmetic) version of the Artin braid group
on \(n+1\) letters, and we can again restate the first assertion of the
theorem by the catch-phrase ``an irreducible local system on \(\mathbb{A}^1\)
minus \(n\) points which is cohomologically rigid extends to a representation
of the braid group on \(n+1\) letters''.

\section{The complex analytic situation}\label{sec:5.8}

\sourceheading[thm:5.8.1]{Theorem 5.8.1}
Over \(\mathbb{C}\), suppose given \(n\geq 2\) an integer,
\(\alpha_1,\alpha_2,\ldots,\alpha_n\) a set of \(n\) distinct points of
\(\mathbb{A}^1(\mathbb{C})\), \(\ell\) a prime number, \(N\geq 1\) an integer,
\(\zeta_N\) a primitive \(N\)'th root of unity in \(\mathbb{C}\), and an
embedding of \(\mathbb{Q}(\zeta_N)\) into
\(\overline{\mathbb{Q}}_\ell\). Let \(\mathcal{F}\) be a cohomologically rigid
object of \(\mathcal{T}_\ell\) which is lisse on
\[
  U:=\mathbb{A}^1-\{\alpha_1,\alpha_2,\ldots,\alpha_n\},
\]
and such that all eigenvalues of all local monodromies of \(\mathcal{F}\) are
\(N\)'th roots of unity. Then there exists on \(U^{\mathrm{an}}\) a local
system \(\mathcal{F}_{\mathrm{cycl}}\) of finite-dimensional
\(\mathbb{Q}(\zeta_N)\)-vector spaces and an isomorphism of
\(\overline{\mathbb{Q}}_\ell\)-local systems on \(U^{\mathrm{an}}\),
\[
  \mathcal{F}^{\mathrm{an}}\cong
  \mathcal{F}_{\mathrm{cycl}}\otimes_{\mathbb{Q}(\zeta_N)}
  \overline{\mathbb{Q}}_\ell.
\]

\begin{proof}
Any lisse, rank one, \(\overline{\mathbb{Q}}_\ell\)-sheaf \(\mathcal{L}\) on
\(U\) whose local monodromies all have order dividing \(N\) is a homomorphism
from \(\pi_1(U)\) to
\(\mu_N(\overline{\mathbb{Q}}_\ell)=\mu_N(\mathbb{Q}(\zeta_N))\). Now
\(\mathcal{L}^{\mathrm{an}}\) is the same homomorphism, restricted to
\(\pi_1(U^{\mathrm{an}})\), so it too has values in
\(\mu_N(\mathbb{Q}(\zeta_N))\), so may be viewed as a rank one
\(\mathbb{Q}(\zeta_N)\)-local system \(\mathcal{L}_{\mathrm{cycl}}\) on
\(U^{\mathrm{an}}\) which sits in an isomorphism
\[
  \mathcal{L}^{\mathrm{an}}\cong
  \mathcal{L}_{\mathrm{cycl}}\otimes_{\mathbb{Q}(\zeta_N)}
  \overline{\mathbb{Q}}_\ell.
\]

We proceed by induction on the generic rank \(r(\mathcal{F})\) of
\(\mathcal{F}\). If \(r(\mathcal{F})=1\), we are done, by the above
discussion. Suppose that \(r(\mathcal{F})\geq 2\). By the
\hyperref[thm:5.2.1]{Main Theorem~5.2.1}, we know that there exists on \(U\) a rank
one lisse \(\mathcal{L}\) of order dividing \(N\), a nontrivial character
\(\chi\) of order dividing \(N\), a lisse cohomologically rigid \(\mathcal{G}\)
with \(r(\mathcal{G})<r(\mathcal{F})\) and with all eigenvalues of all local
monodromies \(N\)'th roots of unity, and an isomorphism on \(U\)
\[
  \mathcal{F}[1]:=
  \mathcal{L}\otimes j^*(j_*\mathcal{G}[1]*_{\mathrm{mid}+}
  j_{0*}\mathcal{L}_\chi[1]).
\]
By induction, the \(\mathbb{Q}(\zeta_N)\)-local system
\(\mathcal{G}_{\mathrm{cycl}}\) exists on \(U^{\mathrm{an}}\), as do
\(\mathcal{L}_{\mathrm{cycl}}\) and \(\mathcal{L}_{\chi,\mathrm{cycl}}\). We
define \(\mathcal{F}_{\mathrm{cycl}}\) on \(U^{\mathrm{an}}\) to be
\[
  \mathcal{F}_{\mathrm{cycl}}[1]:=
  \mathcal{L}_{\mathrm{cycl}}\otimes
  j^*(j_*\mathcal{G}_{\mathrm{cycl}}[1]*_{\mathrm{mid}+}
  j_{0*}\mathcal{L}_{\chi,\mathrm{cycl}}[1]).
\]
By the comparison theorem, \(\mathcal{F}_{\mathrm{cycl}}\) is a local system
of finite-dimensional \(\mathbb{Q}(\zeta_N)\)-vector spaces which sits in an
isomorphism of \(\overline{\mathbb{Q}}_\ell\)-local systems on
\(U^{\mathrm{an}}\),
\[
  \mathcal{F}^{\mathrm{an}}\cong
  \mathcal{F}_{\mathrm{cycl}}\otimes_{\mathbb{Q}(\zeta_N)}
  \overline{\mathbb{Q}}_\ell.
\]
\end{proof}

\sourcepara{5.8.2}
What happens if we drop the quasiunipotence hypothesis? Over \(\mathbb{C}\),
\(I(\infty)\) and each of the groups \(I(\alpha_i)\) are canonically the group
\[
  \varprojlim_N\mu_N(\mathbb{C}),
\]
a group which has a canonical generator, namely
\(\{\exp(2\pi i/N!)\}_N\). Thus we may speak of ``the eigenvalues of local
monodromy'', meaning the eigenvalues of the action of this canonical generator,
of a lisse \(\overline{\mathbb{Q}}_\ell\)-sheaf \(\mathcal{F}\) on
\(\mathbb{A}^1-\{\alpha_1,\alpha_2,\ldots,\alpha_n\}\) at any of the points
\(\infty\) or \(\alpha_i\) where it possibly fails to be lisse.

Repeating the above proof yields:

\sourceheading[thm:5.8.3]{Theorem 5.8.3}
Over \(\mathbb{C}\), suppose given \(n\geq 2\) an integer,
\(\alpha_1,\alpha_2,\ldots,\alpha_n\) a set of \(n\) distinct points of
\(\mathbb{A}^1(\mathbb{C})\), \(\ell\) a prime number, \(\Gamma\) a subgroup
of \(\overline{\mathbb{Q}}_\ell^\times\), \(K\) the subfield
\(\mathbb{Q}(\Gamma)\) of \(\overline{\mathbb{Q}}_\ell\). Let \(\mathcal{F}\)
be a cohomologically rigid object of \(\mathcal{T}_\ell\) which is lisse on
\[
  U:=\mathbb{A}^1-\{\alpha_1,\alpha_2,\ldots,\alpha_n\},
\]
and such that all eigenvalues of all local monodromies of \(\mathcal{F}\) lie
in \(\Gamma\). Then there exists on \(U^{\mathrm{an}}\) a local system
\(\mathcal{F}_{\Gamma}\) of finite-dimensional \(K\)-vector spaces and an
isomorphism of \(\overline{\mathbb{Q}}_\ell\)-local systems on
\(U^{\mathrm{an}}\),
\[
  \mathcal{F}^{\mathrm{an}}\cong
  \mathcal{F}_{\Gamma}\otimes_K\overline{\mathbb{Q}}_\ell.
\]

\section{Return to the original question}\label{sec:5.9}

\sourcepara{5.9.1}
We now return to the question with which we began: what is the structure of
irreducible local systems of finite-dimensional \(\mathbb{C}\)-vector spaces
on \(U^{\mathrm{an}}\) which are physically rigid? We know from
\hyperref[thm:1.1.2]{Theorem~1.1.2} that for such local systems, physical
rigidity is equivalent to cohomological rigidity. Thus it ``suffices'' to
understand the structure of cohomologically rigid irreducible local systems of
finite-dimensional \(\mathbb{C}\)-vector spaces on \(U^{\mathrm{an}}\). Rather
than deal with it directly, we reduce it to the \(\ell\)-adic case on \(U\).
This reduction is made possible by the following standard result, which we
spell out for the convenience of the reader.

\sourceheading[prop:5.9.2]{Proposition 5.9.2}
Over \(\mathbb{C}\), suppose given \(n\geq 2\) an integer,
\(\alpha_1,\alpha_2,\ldots,\alpha_n\) a set of \(n\) distinct points of
\(\mathbb{A}^1(\mathbb{C})\),
\[
  U:=\mathbb{A}^1-\{\alpha_1,\alpha_2,\ldots,\alpha_n\},
\]
and \(\mathcal{F}_{\mathbb{C},\mathrm{an}}\) a local system of
finite-dimensional \(\mathbb{C}\)-vector spaces on \(U^{\mathrm{an}}\). There
exists an integer \(N\geq 1\) such that for all primes \(\ell\) not dividing
\(N\), there exists an isomorphism of fields
\(\iota:\mathbb{C}\cong\overline{\mathbb{Q}}_\ell\), a lisse
\(\overline{\mathbb{Q}}_\ell\)-sheaf \(\mathcal{F}_\ell\) on \(U\), and an
isomorphism of \(\overline{\mathbb{Q}}_\ell\)-local systems on
\(U^{\mathrm{an}}\),
\[
  (\mathcal{F}_\ell)^{\mathrm{an}}\cong
  \mathcal{F}_{\mathbb{C},\mathrm{an}}\otimes_{\mathbb{C}}
  \overline{\mathbb{Q}}_\ell.
\]

\begin{proof}
Let \(n:=\operatorname{rank}(\mathcal{F}_{\mathbb{C},\mathrm{an}})\) be the
rank of \(\mathcal{F}_{\mathbb{C},\mathrm{an}}\). Once we pick a base point in
\(U^{\mathrm{an}}:=U(\mathbb{C})\), we may interpret
\(\mathcal{F}_{\mathbb{C},\mathrm{an}}\) as a homomorphism of groups
\[
  \rho:\pi_1(U^{\mathrm{an}})\to\operatorname{GL}(n,\mathbb{C}).
\]
We know that \(\pi_1(U^{\mathrm{an}})\) is a finitely generated group, and
that its profinite completion is the profinite group \(\pi_1(U)\).

Because \(\pi_1(U^{\mathrm{an}})\) is finitely generated, \(\rho\) takes values
in \(\operatorname{GL}(n,R)\) for some subring \(R\) of \(\mathbb{C}\) which
is finitely generated as a \(\mathbb{Z}\)-algebra. (For instance, if
\(\{\gamma_i\}_i\) is a finite set of generators, we may take for \(R\) the
ring \(\mathbb{Z}[\text{entries of all }\rho(\gamma_i)\text{ and of all }
\rho(\gamma_i^{-1})]\).) Let us admit temporarily the truth of the following
lemma, which is certainly well-known, cf. \cite[6.1.2 (A'')]{BBD}, but for
which I do not know an explicit reference:

\sourceheading[lem:5.9.3]{Lemma 5.9.3}
Let \(R\) be a subring of \(\mathbb{C}\) which is finitely generated as a
\(\mathbb{Z}\)-algebra. Then there exists an integer \(N\geq 1\) such that for
every prime number \(\ell\) which does not divide \(N\), there exists a finite
extension \(E_\lambda\) of \(\mathbb{Q}_\ell\) with integer ring
\(\mathcal{O}_\lambda\), and an isomorphism of fields
\(\iota:\mathbb{C}\cong\overline{\mathbb{Q}}_\ell\) under which
\(\iota(R)\subset\mathcal{O}_\lambda\).

\sourcepara{5.9.4}
Granted the lemma, for any \(\ell\) prime to \(N\) we get that
\(\mathcal{F}_{\mathbb{C},\mathrm{an}}\) has an
\(\mathcal{O}_\lambda\)-form. As the group
\(\operatorname{GL}(n,\mathcal{O}_\lambda)\) is profinite, the composite
homomorphism
\[
  \rho_\lambda:\pi_1(U^{\mathrm{an}})\to
  \operatorname{GL}(n,R)\subset\operatorname{GL}(n,\mathcal{O}_\lambda)
\]
extends to a continuous homomorphism of profinite completions
\[
  \rho_{\lambda,\mathrm{alg}}:\pi_1(U)\to
  \operatorname{GL}(n,\mathcal{O}_\lambda).
\]
Corresponding to \(\rho_{\lambda,\mathrm{alg}}\), we have a lisse
\(\mathcal{O}_\lambda\)-sheaf, say \(\mathcal{F}_\lambda\), and in terms of
\(\mathcal{F}_\lambda\) we define
\[
  \mathcal{F}_\ell:=\mathcal{F}_\lambda
  \otimes_{\mathcal{O}_\lambda}\overline{\mathbb{Q}}_\ell.
\]
\end{proof}

\sourceheading[proofhead:5.9.5]{5.9.5 Proof of \hyperref[lem:5.9.3]{\textcolor{blue}{Lemma~5.9.3}}}
Denote by \(R_{\mathbb{Q}}\) the \(\mathbb{Q}\)-subalgebra of \(\mathbb{C}\)
generated by \(R\). Thus \(R_{\mathbb{Q}}\) is a finitely generated
\(\mathbb{Q}\)-algebra, to which we apply Noether normalization
\cite[2.5]{AK}: there exists a finite collection of elements
\(x_1,\ldots,x_n\) in \(R_{\mathbb{Q}}\) which are algebraically independent
over \(\mathbb{Q}\), and such that \(R_{\mathbb{Q}}\) is integral over
\(\mathbb{Q}[x_1,\ldots,x_n]\). Notice that the ring
\(\mathbb{Q}[x_1,\ldots,x_n]\) is unchanged if we replace each \(x_i\) by a
nonzero integer multiple of itself. Since each \(x_i\) is in
\(R_{\mathbb{Q}}\), after so replacing it we may assume that each \(x_i\) lies
in \(R\).

Now every element \(r\) in \(R\) is integral over
\(\mathbb{Q}[x_1,\ldots,x_n]\). Writing an equation of integral dependence for
\(r\) over \(\mathbb{Q}[x_1,\ldots,x_n]\), we see that for some integer
\(N(r)\geq 1\), \(r\) is integral over
\(\mathbb{Z}[1/N(r)][x_1,\ldots,x_n]\). Because \(R\) is generated as a
\(\mathbb{Z}\)-algebra by finitely many elements \(r_i\) in \(R\), if we
define \(N:=\prod_i N(r_i)\), then each \(r_i\) is integral over
\(\mathbb{Z}[1/N][x_1,\ldots,x_n]\). Since the integral closure of
\(\mathbb{Z}[1/N][x_1,\ldots,x_n]\) in \(R_{\mathbb{Q}}\) is a
\(\mathbb{Z}[1/N]\)-algebra, \(R[1/N]\) is integral over
\(\mathbb{Z}[1/N][x_1,\ldots,x_n]\).

Now pick any prime \(\ell\) not dividing \(N\). Because
\(\mathbb{Q}_\ell\) is uncountable, it has uncountable transcendence degree
over \(\mathbb{Q}\). In particular, there exist \(n\) elements
\(y_1,y_2,\ldots,y_n\) in \(\mathbb{Q}_\ell\) which are algebraically
independent over \(\mathbb{Q}\). Multiplying each \(y_i\) by a power of
\(\ell\), we may further suppose that each \(y_i\) lies in
\(\mathbb{Z}_\ell\).

Using the axiom of choice, there exists an isomorphism
\[
  \iota:\mathbb{C}\cong\overline{\mathbb{Q}}_\ell
\]
such that \(\iota(x_i)=y_i\) for \(i=1,\ldots,n\). Under this isomorphism,
\(\iota(\mathbb{Z}[1/N][x_1,\ldots,x_n])\subset\mathbb{Z}_\ell\), and hence
\(\iota(R[1/N])\) is integral over \(\mathbb{Z}_\ell\). In particular, every
element of \(\iota(R[1/N])\) lies in some finite extension of
\(\mathbb{Q}_\ell\). Since \(R[1/N]\) is finitely generated as a
\(\mathbb{Z}\)-algebra, there exists a single finite extension field
\(E_\lambda\) of \(\mathbb{Q}_\ell\) with
\(\iota(R[1/N])\subset E_\lambda\). As \(\iota(R[1/N])\) is integral over
\(\mathbb{Z}_\ell\), we have
\(\iota(R)\subset\iota(R[1/N])\subset\mathcal{O}_\lambda\).

\section{The category
\texorpdfstring{\(\mathcal{T}_{\mathrm{an}}(U,\Gamma)\)}{T an(U,Gamma)}}\label{sec:5.10}

\sourcepara{5.10.1}
Over \(\mathbb{C}\), suppose given \(n\geq 2\) an integer,
\(\alpha_1,\alpha_2,\ldots,\alpha_n\) a set of \(n\) distinct points of
\(\mathbb{A}^1(\mathbb{C})\),
\[
  U:=\mathbb{A}^1-\{\alpha_1,\alpha_2,\ldots,\alpha_n\},
\]
and \(\Gamma\) a subgroup of \(\mathbb{C}^\times\). We denote by
\(\mathcal{T}_{\mathrm{an}}(U,\Gamma)\) the full subcategory of the category of
all local systems of finite-dimensional \(\mathbb{C}\)-vector spaces on
\(U^{\mathrm{an}}\) whose objects are those which are irreducible, have
nontrivial local monodromy at two or more of the points \(\alpha_i\), and for
which all of eigenvalues of all topological local monodromies lie in
\(\Gamma\).

\sourceheading[thm:5.10.2]{Theorem 5.10.2}
Let \(\chi\) be any nontrivial \(\mathbb{C}^\times\)-valued character with
values in the subgroup \(\Gamma\), and \(\mathcal{L}_{\chi,\mathrm{an}}\) the
corresponding Kummer sheaf on \(\mathbb{G}_{m,\mathrm{an}}\). Then for
\(\mathcal{F}_{\mathrm{an}}\) in \(\mathcal{T}_{\mathrm{an}}(U,\Gamma)\), the
object \(\mathcal{G}_{\mathrm{an}}\) in
\(\mathrm{D}^b_c(U^{\mathrm{an}},\mathbb{C})\) defined by
\[
  \mathcal{G}_{\mathrm{an}}[1]:=
  j^{\mathrm{an}*}(j^{\mathrm{an}}_*\mathcal{F}_{\mathrm{an}}[1]*_{\mathrm{mid}+}
  j_{0,\mathrm{an}*}\mathcal{L}_{\chi,\mathrm{an}}[1]),
\]
lies in \(\mathcal{T}_{\mathrm{an}}(U,\Gamma)\), and its local monodromies are
related to those of \(\mathcal{F}_{\mathrm{an}}\) by the rules of
\hyperref[cor:3.3.6]{Corollary~3.3.6} and
\hyperref[cor:3.3.7]{Corollary~3.3.7}.

\begin{proof}
For any chosen \(\ell\gg 0\), there exists a field isomorphism
\(\iota:\mathbb{C}\cong\overline{\mathbb{Q}}_\ell\) under which both
\(\mathcal{F}_{\mathrm{an}}\) and \(\mathcal{L}_{\chi,\mathrm{an}}\) come from
lisse \(\overline{\mathbb{Q}}_\ell\)-sheaves \(\mathcal{F}\) on \(U\) and
\(\mathcal{L}_\chi\) on \(\mathbb{G}_m\) respectively. If we now view
\(\Gamma\) as a subgroup of \(\overline{\mathbb{Q}}_\ell^\times\), we may
speak of the subgroup \(\Gamma_{\mathrm{values}}\) of the group of continuous
\(\overline{\mathbb{Q}}_\ell^\times\)-valued characters of \(I(0)\) which on
the canonical (we are over \(\mathbb{C}\)) generator take values in
\(\Gamma\). Then \(\mathcal{F}\) lies in
\(\mathcal{T}_\ell(U,\Gamma_{\mathrm{values}})\), and \(\chi\) lies in
\(\Gamma_{\mathrm{values}}\). By
\hyperref[thm:4.3.11]{Theorem~4.3.11}, the object \(\mathcal{G}\) of
\(\mathrm{D}^b_c(U,\overline{\mathbb{Q}}_\ell)\) defined by
\[
  \mathcal{G}[1]:=
  j^*(j_*\mathcal{F}[1]*_{\mathrm{mid}+}j_{0*}\mathcal{L}_\chi[1]),
\]
lies in \(\mathcal{T}_\ell(U,\Gamma_{\mathrm{values}})\), and its local
monodromies are related to those of \(\mathcal{F}\) by the rules of
\hyperref[cor:3.3.6]{Corollary~3.3.6} and
\hyperref[cor:3.3.7]{Corollary~3.3.7}. Passing to \(U^{\mathrm{an}}\) and
applying the inverse of the isomorphism
\(\iota:\mathbb{C}\cong\overline{\mathbb{Q}}_\ell\), we recover our assertions
about \(\mathcal{G}_{\mathrm{an}}\).
\end{proof}

\sourcepara{5.10.3}
We write
\[
  \mathcal{G}_{\mathrm{an}}=\operatorname{MC}_{\chi,\mathrm{an}}
  (\mathcal{F}_{\mathrm{an}}).
\]

\sourcepara{5.10.4}
We can construct a graph whose vertices are the objects of
\(\mathcal{T}_{\mathrm{an}}(U,\Gamma)\), and in which there is an edge joining
two objects \(\mathcal{F}_{\mathrm{an}}\) and \(\mathcal{G}_{\mathrm{an}}\) if
either of the following conditions (a) or (b) holds:
\begin{enumerate}[label=(\alph*)]
\item \(\mathcal{F}_{\mathrm{an}}\) and \(\mathcal{G}_{\mathrm{an}}\) both
have rank \(>1\), and there is a rank one \(\mathbb{C}\)-local system
\(\mathcal{L}_{\mathrm{an}}\) on \(U\), all of whose local monodromy
eigenvalues are in \(\Gamma\), such that
\[
  \mathcal{F}_{\mathrm{an}}\cong
  \mathcal{L}_{\mathrm{an}}\otimes\mathcal{G}_{\mathrm{an}},
  \quad\text{or equivalently}\quad
  \mathcal{G}_{\mathrm{an}}\cong
  \mathcal{L}_{\mathrm{an}}^{-1}\otimes\mathcal{F}_{\mathrm{an}}.
\]

\item there exists a nontrivial character \(\chi\) with values in \(\Gamma\)
such that
\[
  \mathcal{F}_{\mathrm{an}}\cong
  \operatorname{MC}_{\chi,\mathrm{an}}(\mathcal{G}_{\mathrm{an}}),
  \quad\text{or equivalently}\quad
  \mathcal{G}_{\mathrm{an}}\cong
  \operatorname{MC}_{\chi^{-1},\mathrm{an}}(\mathcal{F}_{\mathrm{an}}).
\]
\end{enumerate}

\sourceheading[thm:5.10.5]{Main Theorem 5.10.5 (complex analytic version of \hyperref[thm:5.3.3]{\textcolor{blue}{Main Theorem~5.3.3}})}
If \(\mathcal{F}_{\mathrm{an}}\) is an object of
\(\mathcal{T}_{\mathrm{an}}(U,\Gamma)\) which is rigid, in either of the
equivalent senses of being physically rigid or of being cohomologically rigid,
then \(\mathcal{F}_{\mathrm{an}}\) (as vertex of the above graph) is connected
to an object of rank one in \(\mathcal{T}_{\mathrm{an}}(U,\Gamma)\), and its
distance to such an object is at most
\(2(\operatorname{rank}(\mathcal{F}_{\mathrm{an}})-1)\).

\begin{proof}
Once we have the previous result \hyperref[thm:5.10.2]{Theorem~5.10.2}, we
can repeat, essentially verbatim, the proof already given in the
\(\ell\)-adic case.
\end{proof}

\sourceheading[cor:5.10.6]{Corollary 5.10.6}
In the situation \hyperref[par:5.10.1]{(5.10.1)}, let \(N\geq 2\) be an
integer, and consider the case \(\Gamma=\mu_N(\mathbb{C})\). Thus let
\(\mathcal{F}\) be a local system of finite-dimensional \(\mathbb{C}\)-vector
spaces on \(U^{\mathrm{an}}\) which is irreducible, has nontrivial local
monodromy at two or more of the points \(\alpha_i\), and for which all of
eigenvalues of topological local monodromy are \(N\)'th roots of unity.
Suppose that \(\mathcal{F}\) is rigid. Then
\begin{enumerate}[label=(\arabic*)]
\item \(\mathcal{F}\) has a \(\mathbb{Q}(\zeta_N)\)-form, i.e., for any
embedding of the abstract field \(\mathbb{Q}(\zeta_N)\) into \(\mathbb{C}\),
there exists a local system \(\mathcal{F}_{\mathrm{cycl}}\) of
\(\mathbb{Q}(\zeta_N)\)-vector spaces on \(U^{\mathrm{an}}\) such that
\[
  \mathcal{F}\cong
  \mathcal{F}_{\mathrm{cycl}}\otimes_{\mathbb{Q}(\zeta_N)}\mathbb{C}.
\]

\item Fix a \(\mathbb{Q}(\zeta_N)\)-form \(\mathcal{F}_{\mathrm{cycl}}\) of
\(\mathcal{F}\). For every finite place \(\lambda\) of
\(E:=\mathbb{Q}(\zeta_N)\), with \(\lambda\)-adic completion \(E_\lambda\) and
integer ring \(\mathcal{O}_\lambda\), there exists a lisse
\(E_\lambda\)-sheaf \(\mathcal{F}_\lambda\) on \(U\) and an isomorphism
\[
  (\mathcal{F}_\lambda)^{\mathrm{an}}\cong
  \mathcal{F}_{\mathrm{cycl}}\otimes_E E_\lambda.
\]

\item Fix a \(\mathbb{Q}(\zeta_N)\)-form \(\mathcal{F}_{\mathrm{cycl}}\) of
\(\mathcal{F}\). For every finite place \(\lambda\) of
\(E:=\mathbb{Q}(\zeta_N)\), with \(\lambda\)-adic completion \(E_\lambda\) and
integer ring \(\mathcal{O}_\lambda\), there exists a local system
\(\mathbf{F}_\lambda\) of free \(\mathcal{O}_\lambda\)-modules on \(U\) such
that
\[
  \mathcal{F}_{\mathrm{cycl}}\otimes_E E_\lambda\cong
  (\mathbf{F}_\lambda\otimes_{\mathcal{O}_\lambda}E_\lambda)^{\mathrm{an}}.
\]

\item Fix a \(\mathbb{Q}(\zeta_N)\)-form \(\mathcal{F}_{\mathrm{cycl}}\) of
\(\mathcal{F}\). For every finite place \(\lambda\) of
\(E:=\mathbb{Q}(\zeta_N)\), with \(\lambda\)-adic completion \(E_\lambda\) and
integer ring \(\mathcal{O}_\lambda\), there exists a local system
\(\mathbf{F}^{\mathrm{an}}_\lambda\) of free \(\mathcal{O}_\lambda\)-modules
on \(U^{\mathrm{an}}\) such that
\[
  \mathcal{F}_{\mathrm{cycl}}\otimes_E E_\lambda\cong
  \mathbf{F}^{\mathrm{an}}_\lambda\otimes_{\mathcal{O}_\lambda}E_\lambda.
\]
\end{enumerate}

\begin{proof}
First let us notice that (4) follows trivially from (3): one takes the
\(\mathbf{F}_\lambda\) of (3), and defines \(\mathbf{F}^{\mathrm{an}}_\lambda\)
to be \((\mathbf{F}_\lambda)^{\mathrm{an}}\). Also, (3) follows trivially from
(2), since any lisse \(E_\lambda\)-sheaf \(\mathcal{F}_\lambda\) on \(U\) has
an \(\mathcal{O}_\lambda\)-form.

So it remains to prove (1) and (2). We prove these by induction on the generic
rank of \(\mathcal{F}\). Both are obvious for rank one lisse sheaves
\(\mathcal{L}\) on \(U^{\mathrm{an}}\) all of whose local monodromies are
\(N\)'th roots of unity, and for \(\mathcal{L}_\chi\) with \(\chi\) nontrivial
of order dividing \(N\). Moreover, if (1) and (2) hold for
\(\mathcal{G}_{\mathrm{an}}\), they hold for both
\(\mathcal{L}_{\mathrm{an}}\otimes\mathcal{G}_{\mathrm{an}}\) and
\(\operatorname{MC}_{\chi,\mathrm{an}}(\mathcal{G}_{\mathrm{an}})\) (compare
\hyperref[par:5.5.5.5]{(5.5.5.5)}--\hyperref[par:5.5.5.9]{(5.5.5.9)}). So
by the connectedness properties
\hyperref[thm:5.10.5]{Theorem~5.10.5} of the set of rigid points in the graph
on \(\mathcal{T}_{\mathrm{an}}(U,\Gamma)\), with
\(\Gamma:=\mu_N(\mathbb{C})\), (1) and (2) hold for all \(\mathcal{F}\) in
\(\mathcal{T}_{\mathrm{an}}(U,\Gamma)\).
\end{proof}
